\chapter{ODE linear orde tinggi} \label{ho:chapter}

%%%%%%%%%%%%%%%%%%%%%%%%%%%%%%%%%%%%%%%%%%%%%%%%%%%%%%%%%%%%%%%%%%%%%%%%%%%%%%

\section{ODE linear orde kedua}
\label{solinear:section}

\sectionnotes{1 kuliah, reduksi orde opsional\EPref{,
bagian pertama \S3.1 dalam \cite{EP}}\BDref{,
bagian \S3.1 dan \S3.2 dalam \cite{BD}}}

Perhatikan \emph{\myindex{persamaan diferensial linear orde kedua}} umum
\begin{equation*}
A(x) y'' + B(x)y' + C(x)y = F(x) .
\end{equation*}
Kita sering membagi seluruh persamaan dengan $A(x)$ untuk memperoleh
\begin{equation} \label{sol:eqlin}
y'' + p(x)y' + q(x)y = f(x) ,
\end{equation}
di mana $p(x) = \nicefrac{B(x)}{A(x)}$, $q(x) = \nicefrac{C(x)}{A(x)}$, dan
$f(x) = \nicefrac{F(x)}{A(x)}$.
Kata \emph{linear\index{persamaan linear}} berarti bahwa persamaan tersebut tidak memuat pangkat atau
fungsi dari $y$, $y'$, dan $y''$.

Dalam kasus khusus ketika $f(x) = 0$, kita memperoleh persamaan yang disebut
\emph{homogen\index{persamaan linear homogen}}
\begin{equation} \label{sol:eqlinhom}
y'' + p(x)y' + q(x)y = 0 .
\end{equation}
Kita telah melihat beberapa persamaan homogen linear orde kedua.
\begin{align*}
\qquad y'' + k^2 y & = 0 &
& \text{Dua solusinya adalah:} \quad y_1 = \cos (kx), \quad y_2 = \sin(kx) . \qquad \\
\qquad y'' - k^2 y & = 0 &
& \text{Dua solusinya adalah:} \quad y_1 = e^{kx}, \quad y_2 = e^{-kx} . \qquad
\end{align*}

Jika kita mengetahui dua solusi suatu persamaan linear homogen, kita mengetahui jauh lebih banyak solusi lainnya.

\begin{theorem}[Superposisi]\index{superposisi}
Andaikan $y_1$ dan $y_2$ adalah dua solusi persamaan homogen \eqref{sol:eqlinhom}. Maka
\begin{equation*}
y(x) = C_1 y_1(x) + C_2 y_2(x) ,
\end{equation*}
juga menyelesaikan \eqref{sol:eqlinhom} untuk konstanta sebarang $C_1$ dan $C_2$.
\end{theorem}

Artinya, kita dapat menjumlahkan solusi dan mengalikannya dengan konstanta untuk memperoleh solusi baru yang berbeda. Kita menyebut ekspresi $C_1 y_1 + C_2 y_2$ sebagai \emph{\myindex{kombinasi linear}} dari $y_1$ dan $y_2$. Mari kita buktikan teorema ini; buktinya sangat mencerahkan dan menggambarkan cara kerja persamaan linear.

\medskip

\emph{Bukti:}
Misalkan $y = C_1 y_1 + C_2 y_2$. Maka
\begin{equation*}
\begin{split}
y'' + py' + qy & =
(C_1 y_1 + C_2 y_2)'' + p(C_1 y_1 + C_2 y_2)' + q(C_1 y_1 + C_2 y_2) \\
& = C_1 y_1'' + C_2 y_2'' + C_1 p y_1' + C_2 p y_2' + C_1 q y_1 + C_2 q y_2 \\
& = C_1 ( y_1'' + p y_1' + q y_1 ) + C_2 ( y_2'' + p y_2' + q y_2 ) \\
& = C_1 \cdot 0 + C_2 \cdot 0 = 0 . \qed
\end{split}
\end{equation*}

\medskip

Bukti tersebut bahkan lebih sederhana dinyatakan jika kita menggunakan notasi operator. Sebuah \emph{\myindex{operator}} adalah objek yang menerima fungsi dan menghasilkan fungsi (mirip dengan fungsi, tetapi fungsi menerima bilangan dan menghasilkan bilangan). Definisikan operator $L$ dengan
\begin{equation*}
Ly = y'' + py' + qy .
\end{equation*}
Persamaan diferensial kini menjadi $Ly=0$.
Bahwa operator (dan persamaan) $L$ bersifat \emph{linear}\index{operator linear} berarti $L(C_1y_1 + C_2y_2) = C_1 Ly_1 + C_2 Ly_2$. Seolah-olah kita sedang \myquote{mengalikan} dengan $L$. Bukti di atas menjadi
\begin{equation*}
Ly = L(C_1y_1 + C_2y_2) = C_1 Ly_1 + C_2 Ly_2 = C_1 \cdot 0 + C_2 \cdot 0 = 0 .
\end{equation*}

\medskip

Dua solusi lain untuk persamaan kedua $y'' - k^2y = 0$ adalah $y_1 = \cosh (kx)$ dan $y_2 = \sinh (kx)$. Dengan mengingat definisi $\sinh$ dan $\cosh$, kita mencatat bahwa keduanya merupakan solusi berdasarkan superposisi karena merupakan kombinasi linear dari dua solusi eksponensial:
$\cosh(kx) = \frac{e^{kx}  + e^{-kx}}{2} = (\nicefrac{1}{2}) e^{kx}  +
(\nicefrac{1}{2}) e^{-kx}$ dan
$\sinh(kx) = \frac{e^{kx} - e^{-kx}}{2} = (\nicefrac{1}{2}) e^{kx}  +
(\nicefrac{-1}{2}) e^{-kx}$.

Fungsi $\sinh$ dan $\cosh$ kadang lebih praktis digunakan daripada fungsi eksponensial. Mari kita tinjau beberapa sifatnya:
\begin{align*}
& \cosh 0  = 1 , &   & \sinh 0 = 0 , \\
& \frac{d}{dx} \Bigl[ \cosh x \Bigr] = \sinh x , &  & \frac{d}{dx} \Bigl[ \sinh x \Bigr] = \cosh x , \\
& \cosh^2 x - \sinh^2 x = 1 .
\end{align*}

\begin{exercise}
Turunkan sifat-sifat ini menggunakan definisi $\sinh$ dan $\cosh$ dalam bentuk eksponensial.
\end{exercise}

Persamaan linear memberikan jawaban yang baik dan sederhana bagi pertanyaan keberadaan dan ketunggalan.

\begin{theorem}[Keberadaan dan ketunggalan]\index{keberadaan dan ketunggalan}
Andaikan $p, q, f$ adalah fungsi kontinu pada suatu interval $I$, $a$ adalah suatu bilangan dalam $I$, dan $b_0, b_1$ adalah konstanta. Maka persamaan
\begin{equation*}
y'' + p(x) y' + q(x) y = f(x) ,
\end{equation*}
memiliki tepat satu solusi $y(x)$ yang terdefinisi pada interval $I$ dan memenuhi syarat awal
\begin{equation*}
y(a) = b_0 , \qquad y'(a) = b_1 .
\end{equation*}
\end{theorem}

Sebagai contoh, persamaan $y'' + k^2 y = 0$ dengan $y(0) = b_0$ dan $y'(0) = b_1$ memiliki solusi
\begin{equation*}
y(x) = b_0 \cos (kx) + \frac{b_1}{k} \sin (kx) .
\end{equation*}
Persamaan $y'' - k^2 y = 0$ dengan $y(0) = b_0$ dan $y'(0) = b_1$ memiliki solusi
\begin{equation*}
y(x) = b_0 \cosh (kx) + \frac{b_1}{k} \sinh (kx) .
\end{equation*}
Penggunaan $\cosh$ dan $\sinh$ dalam solusi ini memungkinkan kita menyelesaikan syarat awal dengan lebih rapi daripada jika kita menggunakan eksponensial.

\medskip

Syarat awal untuk ODE orde kedua terdiri atas dua persamaan. Akal sehat memberi tahu kita bahwa jika terdapat dua konstanta sebarang dan dua persamaan, kita semestinya dapat menentukan konstanta-konstanta tersebut dan menemukan solusi persamaan diferensial yang memenuhi syarat awal.

\emph{Pertanyaan:} Andaikan kita menemukan dua solusi berbeda $y_1$ dan $y_2$ untuk persamaan homogen \eqref{sol:eqlinhom}. Dapatkah setiap solusi ditulis (dengan superposisi) dalam bentuk $y = C_1 y_1 + C_2 y_2$?

Jawabannya ya! Asalkan $y_1$ dan $y_2$ cukup berbeda dalam arti berikut. Kita katakan $y_1$ dan $y_2$ \emph{\myindex{bebas linear}} jika yang satu bukan kelipatan konstan dari yang lain.

\begin{theorem}
Misalkan $p, q$ fungsi kontinu. Misalkan $y_1$ dan $y_2$ dua solusi bebas linear dari persamaan homogen \eqref{sol:eqlinhom}. Maka setiap solusi lainnya berbentuk
\begin{equation*}
y = C_1 y_1 + C_2 y_2 .
\end{equation*}
Artinya, $y = C_1 y_1 + C_2 y_2$ adalah solusi umum.
\end{theorem}

Sebagai contoh, kita menemukan solusi $y_1 = \cos x$ dan $y_2 = \sin x$ untuk persamaan $y'' + y = 0$. Mudah dilihat bahwa sinus dan kosinus bukan kelipatan konstan satu sama lain. Memang, jika $\sin x = A \cos x$ untuk suatu konstanta $A$, substitusi $x=0$ akan mengakibatkan $A = 0$. Namun kemudian $\sin x = 0$ untuk semua $x$, yang mustahil. Jadi $y_1$ dan $y_2$ bebas linear. Oleh karena itu,
\begin{equation*}
y = C_1 \cos x + C_2 \sin x
\end{equation*}
adalah solusi umum untuk $y'' + y = 0$.

Untuk dua fungsi, pemeriksaan kebebasan linear cukup sederhana. Berikut contoh lain. Perhatikan $y''-2x^{-2}y = 0$. Maka $y_1 = x^2$ dan $y_2 = \nicefrac{1}{x}$ adalah solusi. Untuk melihat bahwa keduanya bebas linear, andaikan salah satunya merupakan kelipatan yang lain: $y_1 = A y_2$; kita hanya perlu menunjukkan bahwa $A$ tidak mungkin konstan. Dalam kasus ini, $A = \nicefrac{y_1}{y_2} = x^3$, yang jelas bukan konstanta. Jadi $y = C_1 x^2 + C_2 \nicefrac{1}{x}$ adalah solusi umum.

\medskip

Jika Anda memiliki satu solusi taknol untuk suatu persamaan homogen linear orde kedua, Anda dapat menemukan solusi lain. Inilah \emph{\myindex{metode reduksi orde}}. Gagasannya ialah bahwa jika kita entah bagaimana menemukan $y_1$ sebagai solusi $y'' + p(x) y' + q(x) y = 0$, kita mencoba solusi kedua berbentuk $y_2(x) = y_1(x) v(x)$.
Kita hanya perlu menemukan $v$. Kita substitusikan $y_2$ ke persamaan:
\begin{equation*}
\begin{split}
0 =
y_2'' + p(x) y_2' + q(x) y_2
& =
\underbrace{y_1'' v + 2 y_1' v' + y_1 v''}_{y_2''}
+ p(x) \underbrace{( y_1' v + y_1 v' )}_{y_2'}
+ q(x) \underbrace{y_1 v}_{y_2}
\\
& =
y_1 v''
+ (2 y_1' + p(x) y_1) v'
+
\cancelto{0}{\bigl( y_1'' + p(x) y_1' + q(x) y_1 \bigr)} v .
\end{split}
\end{equation*}
Dengan kata lain, $y_1 v'' + (2 y_1' + p(x) y_1) v' = 0$. Dengan menggunakan $w = v'$, kita memperoleh persamaan linear orde pertama $y_1 w' + (2 y_1' + p(x) y_1) w = 0$. Setelah menyelesaikan persamaan ini untuk $w$ (faktor integrasi), kita menemukan $v$ dengan menganti-diferensialkan $w$. Kemudian kita membentuk $y_2$ dengan menghitung $y_1v$. Sebagai contoh, andaikan kita entah bagaimana mengetahui bahwa $y_1 = x$ adalah solusi $y''+x^{-1}y'-x^{-2} y=0$.
Persamaan untuk $w$ kemudian adalah $xw' + 3 w = 0$. Kita menemukan solusi $w = Cx^{-3}$ dan antiturunan $v = \frac{-C}{2x^2}$. Jadi $y_2 = y_1 v = \frac{-C}{2x}$. Sebarang $C$ dapat digunakan, sehingga $C=-2$ menghasilkan $y_2 = \nicefrac{1}{x}$. Dengan demikian, solusi umumnya adalah $y = C_1 x + C_2\nicefrac{1}{x}$.

Karena kita memiliki rumus solusi persamaan linear orde pertama, kita dapat menuliskan rumus untuk $y_2$:
\begin{equation*}
y_2(x) = y_1(x) \int \frac{e^{-\int p(x)\,dx}}{{\bigl(y_1(x)\bigr)}^2} \,dx
\end{equation*}
Namun, jauh lebih mudah mengingat bahwa kita hanya perlu mencoba $y_2(x) = y_1(x) v(x)$ dan menemukan $v(x)$ seperti di atas. Teknik ini juga berlaku untuk persamaan orde lebih tinggi: orde berkurang satu untuk setiap solusi yang ditemukan. Jadi lebih baik mengingat caranya daripada suatu rumus khusus.

\medskip

Kita akan mempelajari solusi persamaan nonhomogen dalam \sectionref{sec:nonhom}. Untuk saat ini kita berfokus mencari solusi umum persamaan homogen.

\subsection{Latihan}

\begin{exercise}
Tunjukkan bahwa $y=e^x$ dan $y=e^{2x}$ bebas linear.
\end{exercise}

\begin{exercise}
Ambil $y'' + 5 y = 10 x + 5$. Temukan (tebak!) sebuah solusi.
\end{exercise}

\begin{exercise}
Buktikan prinsip superposisi untuk persamaan nonhomogen. Andaikan $y_1$ adalah solusi $L y_1 = f(x)$ dan $y_2$ adalah solusi $L y_2 = g(x)$ (operator linear $L$ yang sama). Tunjukkan bahwa $y = y_1+y_2$ menyelesaikan $Ly = f(x) + g(x)$.
\end{exercise}

\begin{exercise}
Untuk persamaan $x^2 y'' - x y' = 0$, temukan dua solusi, tunjukkan bahwa keduanya bebas linear, dan temukan solusi umumnya.
Petunjuk: Coba $y = x^r$.
\end{exercise}

\pagebreak[2]
Persamaan berbentuk $a x^2 y'' + b x y' + c y = 0$ disebut
\emph{persamaan Euler\index{persamaan Euler}} atau
\emph{persamaan Cauchy--Euler\index{persamaan Cauchy--Euler}}.
Persamaan tersebut diselesaikan dengan mencoba $y=x^r$ dan menentukan $r$ (andaikan $x \geq 0$ demi kesederhanaan).

\begin{exercise} \label{sol:eulerex}
\pagebreak[2]
Andaikan ${(b-a)}^2-4ac > 0$.
\begin{tasks}
\task Temukan rumus solusi umum persamaan Euler (lihat di atas) $a x^2 y'' + b x y' + c y = 0$.
Petunjuk: Coba $y=x^r$ dan temukan rumus untuk $r$.
\task Apa yang terjadi ketika ${(b-a)}^2-4ac = 0$ atau ${(b-a)}^2-4ac < 0$?
\end{tasks}
\end{exercise}

Kita akan meninjau kembali kasus ${(b-a)}^2-4ac < 0$ nanti.

\begin{exercise} \label{sol:eulerexln}
Persamaan yang sama seperti dalam \exerciseref{sol:eulerex}.
Andaikan ${(b-a)}^2-4ac = 0$. Temukan rumus solusi umum $a x^2 y'' + b x y' + c y = 0$. Petunjuk: Coba $y=x^r \ln x$ untuk solusi kedua.
\end{exercise}

\begin{exercise}[reduksi orde] \label{exercise:reductionoforder}
Andaikan $y_1$ adalah solusi $y'' + p(x) y' + q(x) y = 0$. Dengan menyubstitusikannya langsung ke persamaan, tunjukkan bahwa
\begin{equation*}
y_2(x) = y_1(x) \int \frac{e^{-\int p(x)\,dx}}{{\bigl(y_1(x)\bigr)}^2} \,dx
\end{equation*}
juga merupakan solusi.
\end{exercise}

\begin{exercise}[\myindex{Persamaan Chebyshev orde 1}]
Ambil $(1-x^2)y''-xy' + y = 0$.
\begin{tasks}
\task Tunjukkan bahwa $y=x$ adalah solusi.
\task Gunakan reduksi orde untuk menemukan solusi kedua yang bebas linear.
\task Tuliskan solusi umumnya.
\end{tasks}
\end{exercise}

\begin{exercise}[\myindex{Persamaan Hermite orde 2}]
Ambil $y''-2xy' + 4y = 0$.
\begin{tasks}
\task Tunjukkan bahwa $y=1-2x^2$ adalah solusi.
\task Gunakan reduksi orde untuk menemukan solusi kedua yang bebas linear.
(Integral tentu boleh dibiarkan dalam rumus.)
\task Tuliskan solusi umumnya.
\end{tasks}
\end{exercise}

\setcounter{exercise}{100}

\begin{exercise}
Apakah $\sin(x)$ dan $e^x$ bebas linear? Berikan alasan.
\end{exercise}
\exsol{%
Ya. Untuk membenarkannya, cobalah mencari konstanta $A$ sedemikian sehingga $\sin(x) = A e^x$ untuk semua $x$.
}

\begin{exercise}
Apakah $e^x$ dan $e^{x+2}$ bebas linear? Berikan alasan.
\end{exercise}
\exsol{%
Tidak. $e^{x+2} = e^2 e^x$.
}

\begin{exercise}
Tebak sebuah solusi untuk $y'' + y' + y= 5$.
\end{exercise}
\exsol{%
$y=5$
}

\begin{exercise}
Temukan solusi umum $x y'' + y' = 0$. Petunjuk: Ini adalah ODE orde pertama dalam $y'$.
\end{exercise}
\exsol{%
$y=C_1 \ln(x) + C_2$
}

\begin{exercise}
Tuliskan sebuah persamaan (tebak) yang mempunyai solusi $e^x$ dan $e^{2x}$. Petunjuk: Coba persamaan berbentuk $y''+Ay'+By = 0$ untuk konstanta $A$ dan $B$, substitusikan $e^x$ dan $e^{2x}$, lalu tentukan $A$ dan $B$.
\end{exercise}
\exsol{%
$y''-3y'+2y = 0$
}

%%%%%%%%%%%%%%%%%%%%%%%%%%%%%%%%%%%%%%%%%%%%%%%%%%%%%%%%%%%%%%%%%%%%%%%%%%%%%%

\sectionnewpage
\section{ODE linear orde kedua dengan koefisien konstan}
\label{sec:ccsol}

%mbxINTROSUBSECTION

\sectionnotes{lebih dari 1 kuliah\EPref{,
bagian kedua \S3.1 dalam \cite{EP}}\BDref{,
\S3.1 dalam \cite{BD}}}

\subsection{Menyelesaikan persamaan berkoefisien konstan}

Perhatikan masalah
\begin{equation*}
y''-6y'+8y = 0, \qquad y(0) = - 2, \qquad y'(0) = 6 .
\end{equation*}
Ini adalah persamaan homogen linear orde kedua dengan koefisien konstan. \emph{Koefisien konstan\index{koefisien konstan}} berarti fungsi-fungsi di depan $y''$, $y'$, dan $y$ adalah konstanta; semuanya tidak bergantung pada $x$.

Untuk menebak solusi, pikirkan suatu fungsi yang pada dasarnya tetap sama ketika didiferensialkan, sehingga kita dapat mengambil fungsi itu beserta turunannya, menjumlahkan beberapa kelipatannya, dan memperoleh nol. Ya, yang dimaksud adalah fungsi eksponensial.

Kita mencoba\footnote{%
Membuat tebakan terdidik dengan beberapa parameter yang perlu ditentukan merupakan teknik yang begitu penting dalam persamaan diferensial sehingga orang kadang menggunakan istilah mewah untuk tebakan semacam itu: \emph{\myindex{ansatz}}, bahasa Jerman untuk \myquote{penempatan awal suatu alat pada benda kerja.} Ya, orang Jerman mempunyai kata untuk itu.}
solusi berbentuk $y = e^{rx}$. Maka $y' = r e^{rx}$ dan $y'' = r^2 e^{rx}$. Substitusikan untuk memperoleh
\begin{align*}
y''-6y'+8y & = 0 , \\
\underbrace{r^2 e^{rx}}_{y''} -6 \underbrace{r e^{rx}}_{y'}+8 \underbrace{e^{rx}}_{y} & = 0 , \\
r^2 -6 r +8 & = 0 \qquad \text{(bagi seluruhnya dengan } e^{rx} \text{)},\\
(r-2)(r-4) & = 0 .
\end{align*}
Jadi, jika $r=2$ atau $r=4$, maka $e^{rx}$ adalah solusi. Ambil $y_1 = e^{2x}$ dan $y_2 = e^{4x}$.

\begin{exercise}
Periksa bahwa $y_1$ dan $y_2$ adalah solusi.
\end{exercise}

Fungsi $e^{2x}$ dan $e^{4x}$ bebas linear. Jika tidak, kita dapat menulis $e^{4x} = C e^{2x}$ untuk suatu konstanta $C$, yang mengakibatkan $e^{2x} = C$ untuk semua $x$, dan ini jelas mustahil. Jadi solusi umumnya dapat ditulis sebagai
\begin{equation*}
y = C_1 e^{2x} + C_2 e^{4x} .
\end{equation*}
Kita perlu menentukan $C_1$ dan $C_2$. Untuk menerapkan syarat awal, mula-mula temukan $y' = 2 C_1 e^{2x} + 4 C_2 e^{4x}$. Substitusikan $x=0$ ke $y$ dan $y'$, lalu selesaikan.
\begin{equation*}
\begin{aligned}
-2 & = y(0) = C_1 + C_2 , \\
6 & = y'(0) = 2 C_1 + 4 C_2 .
\end{aligned}
\end{equation*}
Gunakan aljabar matriks atau cukup aljabar sekolah menengah. Misalnya, bagi persamaan kedua dengan 2 untuk memperoleh $3 = C_1 + 2 C_2$, lalu kurangkan kedua persamaan sehingga diperoleh $5 = C_2$. Kemudian $C_1 = -7$ karena $-2 = C_1 + 5$. Jadi solusi yang dicari adalah
\begin{equation*}
y = -7 e^{2x} + 5 e^{4x} .
\end{equation*}

\medskip

Kita perumum contoh ini menjadi suatu metode. Andaikan kita mempunyai persamaan
\begin{equation} \label{ccsol:eq}
a y'' + b y' + c y = 0 ,
\end{equation}
di mana $a, b, c$ adalah konstanta. Coba solusi $y = e^{rx}$ untuk memperoleh
\begin{equation*}
a r^2 e^{rx} + b r e^{rx} + c e^{rx} = 0 .
\end{equation*}
Bagi dengan $e^{rx}$ untuk memperoleh apa yang disebut \emph{\myindex{persamaan karakteristik}} ODE tersebut:
\begin{equation*}
a r^2 + b r + c = 0 .
\end{equation*}
Tentukan $r$ dengan menggunakan \myindex{rumus kuadrat}:
\begin{equation*}
r_1, r_2 = \frac{-b \pm \sqrt{b^2 - 4ac}}{2a} .
\end{equation*}
Untuk saat ini, andaikan $b^2 -4ac \geq 0$ sehingga $r_1$ dan $r_2$ real. Maka $e^{r_1 x}$ dan $e^{r_2 x}$ adalah solusi. Masih ada kesulitan jika $r_1 = r_2$, tetapi mudah diatasi.

\begin{theorem}
Andaikan $r_1$ dan $r_2$ adalah akar persamaan karakteristik.
\begin{enumerate}[(i)]
\item Jika $r_1$ dan $r_2$ berbeda dan real (ketika $b^2 - 4ac > 0$), maka \eqref{ccsol:eq} memiliki solusi umum
\begin{equation*}
\mybxbg{~~
y = C_1 e^{r_1 x} + C_2 e^{r_2 x} .
~~}
\end{equation*}
\item Jika $r_1 = r_2$ (terjadi ketika $b^2 - 4ac = 0$), maka \eqref{ccsol:eq} memiliki solusi umum
\begin{equation*}
\mybxbg{~~
y = (C_1 + C_2 x)\, e^{r_1 x} .
~~}
\end{equation*}
\end{enumerate}
\end{theorem}

\begin{example} \label{example:expsecondorder}
Selesaikan
\begin{equation*}
y'' - k^2 y = 0 .
\end{equation*}
Persamaan karakteristiknya adalah $r^2 - k^2 = 0$ atau $(r-k)(r+k) = 0$. Akibatnya, $e^{-k x}$ dan $e^{kx}$ adalah dua solusi bebas linear, dan solusi umumnya
\begin{equation*}
y = C_1 e^{kx} + C_2e^{-kx} .
\end{equation*}
Karena $\cosh s = \frac{e^s+e^{-s}}{2}$ dan $\sinh s = \frac{e^s-e^{-s}}{2}$, solusi umum juga dapat ditulis sebagai
\begin{equation*}
y = D_1 \cosh(kx) + D_2 \sinh(kx) .
\end{equation*}
\end{example}

\begin{example}
Temukan solusi umum
\begin{equation*}
y'' -8 y' + 16 y = 0 .
\end{equation*}
Persamaan karakteristiknya adalah $r^2 - 8 r + 16 = {(r-4)}^2 = 0$. Persamaan ini mempunyai akar ganda $r_1 = r_2 = 4$. Jadi solusi umumnya
\begin{equation*}
y = (C_1 + C_2 x)\, e^{4 x} = C_1 e^{4x} + C_2 x e^{4x} .
\end{equation*}

\begin{exercise}
Periksa bahwa $e^{4x}$ dan $x e^{4x}$ bebas linear.
\end{exercise}

Ada baiknya memeriksa pekerjaan Anda. Bahwa $e^{4x}$ menyelesaikan persamaan sudah jelas. Mari periksa bahwa $x e^{4x}$ menyelesaikannya. Hitung $y' = e^{4x} + 4xe^{4x}$ dan $y'' = 8 e^{4x} + 16xe^{4x}$. Substitusikan,
\begin{equation*}
y'' - 8 y' + 16 y =
8 e^{4x} + 16xe^{4x} - 8(e^{4x} + 4xe^{4x}) + 16 xe^{4x} =
0 .
\end{equation*}
\end{example}

Dalam arti tertentu, akar ganda jarang terjadi. Jika koefisien dipilih secara acak, akar ganda kecil kemungkinannya. Namun, ada beberapa masalah nyata tempat akar ganda muncul secara alami (misalnya sistem massa-pegas teredam kritis, seperti yang akan kita lihat).

Berikut argumen singkat mengapa solusi $x e^{r x}$ berlaku untuk akar ganda. Akar ganda adalah kasus limit dari dua akar berbeda yang sangat berdekatan. Perhatikan bahwa $\frac{e^{r_2 x} - e^{r_1 x}}{r_2 - r_1}$ adalah solusi ketika akarnya berbeda. Ketika mengambil limit saat $r_1$ menuju $r_2$, sebenarnya kita mengambil turunan $e^{rx}$ dengan $r$ sebagai variabel. Jadi limitnya adalah $x e^{rx}$, dan dengan demikian ini merupakan solusi dalam kasus akar ganda. Kita mencatat bahwa dalam beberapa perhitungan numerik, dua akar yang sangat berdekatan dapat menimbulkan ketidakstabilan numerik, sedangkan akar ganda tidak.

\subsection{Bilangan kompleks dan rumus Euler}

Suatu polinomial dapat mempunyai akar kompleks. Persamaan $r^2 + 1 = 0$ tidak mempunyai akar real, tetapi mempunyai dua akar kompleks. Di sini kita meninjau beberapa sifat bilangan kompleks\index{bilangan kompleks}.

Bilangan kompleks mungkin tampak sebagai konsep aneh, barangkali karena istilahnya. Tidak ada yang imajiner atau rumit tentang bilangan kompleks. Bilangan kompleks hanyalah pasangan bilangan real, $(a,b)$, yaitu sebuah titik pada bidang. Kita menjumlahkan bilangan kompleks dengan cara langsung: $(a,b)+(c,d)=(a+c,b+d)$. Kita mendefinisikan perkalian\index{perkalian bilangan kompleks} dengan
\begin{equation*}
(a,b) \times (c,d) \overset{\text{def}}{=} (ac-bd,ad+bc) .
\end{equation*}
Ternyata, dengan aturan perkalian ini, semua sifat aritmetika standar berlaku. Lebih lanjut, dan yang terpenting, $(0,1) \times (0,1) = (-1,0)$.

Umumnya kita menulis $(a,b)$ sebagai $a+ib$, dan memperlakukan $i$ seolah-olah suatu peubah. Ketika $b$ nol, $(a,0)$ hanyalah bilangan $a$. Kita melakukan aritmetika dengan bilangan kompleks sebagaimana dengan polinomial. Sifat yang baru disebutkan menjadi $i^2 = -1$. Jadi setiap kali melihat $i^2$, kita menggantinya dengan $-1$. Sebagai contoh,
\begin{equation*}
(2+3i)(4i) - 5i =
(2\times 4)i + (3 \times 4) i^2 - 5i
= 8i + 12 (-1) - 5i
= -12 + 3i .
\end{equation*}

Bilangan $i$ dan $-i$ adalah dua akar $r^2 + 1 = 0$. Sebagian insinyur menggunakan huruf $j$ alih-alih $i$ untuk akar kuadrat $-1$. Kita mengikuti konvensi matematikawan dan menggunakan $i$.

\begin{exercise}
Pastikan Anda memahami (dan dapat membenarkan) identitas berikut:
\begin{tasks}(2)
\task $i^2 = -1$, $i^3 = -i$, $i^4 = 1$,
\task $\frac{1}{i} = -i$,
\task $(3-7i)(-2-9i) = \cdots = -69-13i$,
\task $(3-2i)(3+2i) = 3^2 - {(2i)}^2 = 3^2 + 2^2 = 13$,
\task $\frac{1}{3-2i} = \frac{1}{3-2i} \frac{3+2i}{3+2i} = \frac{3+2i}{13}
= \frac{3}{13}+\frac{2}{13}i$.
\end{tasks}
\end{exercise}

\pagebreak[2]
Kita juga mendefinisikan eksponensial $e^{a+ib}$ dari bilangan kompleks dengan menuliskan deret Taylor lalu menyubstitusikan bilangan kompleks tersebut. Sebagian besar sifat eksponensial mengikuti dari deret Taylornya, sehingga sifat tersebut juga berlaku untuk eksponensial kompleks. Sebagai contoh, $e^{x+y} = e^x e^y$. Khususnya, $e^{a+ib} = e^a e^{ib}$. Jadi jika kita dapat menghitung $e^{ib}$, kita dapat menghitung $e^{a+ib}$. Untuk $e^{ib}$, kita menggunakan apa yang disebut \emph{\myindex{rumus Euler}}.

\begin{theorem}[Rumus Euler] \label{eulersformula}
\begin{equation*}
\mybxbg{~~
e^{i \theta} = \cos \theta + i \sin \theta
\qquad \text{ dan } \qquad
e^{- i \theta} = \cos \theta - i \sin \theta .
~~}
\end{equation*}
\end{theorem}

Dengan kata lain, $e^{a+ib} = e^a \bigl( \cos(b) + i \sin(b) \bigr) = e^a \cos(b) + i e^a \sin(b)$.

\begin{exercise}
Dengan menggunakan rumus Euler, periksa identitas:
\begin{equation*}
\cos \theta = \frac{e^{i \theta} + e^{-i \theta}}{2}
\qquad \text{dan} \qquad
\sin \theta = \frac{e^{i \theta} - e^{-i \theta}}{2i}.
\end{equation*}
\end{exercise}

\begin{exercise}[Identitas sudut ganda]
Mulailah dengan $e^{i(2\theta)} = {\bigl(e^{i \theta} \bigr)}^2$. Gunakan rumus Euler pada setiap ruas untuk menyimpulkan:
\begin{equation*}
\cos (2\theta) = \cos^2 \theta - \sin^2 \theta
\qquad \text{dan} \qquad
\sin (2\theta) = 2 \sin \theta \cos \theta .
\end{equation*}
\end{exercise}

Untuk bilangan kompleks $a+ib$, kita menyebut $a$ sebagai \emph{\myindex{bagian real}} dan $b$ sebagai \emph{\myindex{bagian imajiner}} bilangan tersebut. Notasi berikut sering digunakan:
\begin{equation*}
\operatorname{Re}(a+ib) = a
\qquad \text{dan} \qquad
\operatorname{Im}(a+ib) = b.
\end{equation*}

\subsection{Akar kompleks}

Andaikan persamaan diferensial $ay'' + by' + cy = 0$ mempunyai persamaan karakteristik $a r^2 + b r + c = 0$ dengan \myindex{akar kompleks}. Berdasarkan rumus kuadrat, akarnya adalah $\frac{-b \pm \sqrt{b^2 - 4ac}}{2a}$. Akar-akar ini kompleks jika $b^2 - 4ac < 0$. Dalam kasus ini, kita menulis akarnya sebagai
\begin{equation*}
r_1, r_2 = \frac{-b}{2a} \pm i\frac{\sqrt{4ac - b^2}}{2a} .
\end{equation*}
Seperti terlihat, kita memperoleh sepasang akar berbentuk $\alpha \pm i \beta$. Kita masih dapat menulis solusi sebagai
\begin{equation*}
y = C_1 e^{(\alpha+i\beta)x} + C_2 e^{(\alpha-i\beta)x} .
\end{equation*}
Namun, eksponensialnya kini bernilai kompleks. Kita perlu mengizinkan $C_1$ dan $C_2$ berupa bilangan kompleks untuk memperoleh solusi bernilai real (yang kita cari). Meskipun pendekatan ini tidak bermasalah, perhitungannya dapat menjadi lebih sulit dan umumnya lebih disukai mencari dua solusi bernilai real.

\hyperref[eulersformula]{Rumus Euler} datang menolong. Misalkan
\begin{equation*}
y_1 = e^{(\alpha+i\beta)x} \qquad \text{dan} \qquad y_2 = e^{(\alpha-i\beta)x} .
\end{equation*}
Maka
\begin{equation*}
\begin{aligned}
y_1 & = e^{\alpha x} \cos (\beta x) + i e^{\alpha x} \sin (\beta x) , \\
y_2 & = e^{\alpha x} \cos (\beta x) - i e^{\alpha x} \sin (\beta x) .
\end{aligned}
\end{equation*}
Kombinasi linear solusi juga merupakan solusi. Jadi,
\begin{equation*}
\begin{aligned}
y_3 & = \frac{y_1 + y_2}{2} = e^{\alpha x} \cos (\beta x) , \\
y_4 & = \frac{y_1 - y_2}{2i} = e^{\alpha x} \sin (\beta x) ,
\end{aligned}
\end{equation*}
juga merupakan solusi. Mudah dilihat bahwa $y_3$ dan $y_4$ bebas linear (bukan kelipatan satu sama lain). Jadi solusi umum dapat ditulis dengan $y_3$ dan $y_4$. Karena keduanya bernilai real, konstanta sebarang dalam solusi umum tidak perlu berupa bilangan kompleks. Kita rangkum hasil ini sebagai teorema.

\begin{theorem}
Ambil persamaan
\begin{equation*}
ay'' + by' + cy = 0 .
\end{equation*}
Jika persamaan karakteristik mempunyai akar $\alpha \pm i \beta$ (ketika $b^2 - 4ac < 0$), maka solusi umumnya
\begin{equation*}
\mybxbg{~~
y = C_1 e^{\alpha x} \cos (\beta x) + C_2 e^{\alpha x} \sin (\beta x) .
~~}
\end{equation*}
\end{theorem}

\begin{example} \label{example:sincossecondorder}
Temukan solusi umum $y'' + k^2 y = 0$, untuk konstanta $k > 0$.

Persamaan karakteristiknya adalah $r^2 + k^2 = 0$, dan akarnya $r = \pm ik$. Berdasarkan teorema, solusi umumnya
\begin{equation*}
y = C_1 \cos (kx) + C_2 \sin (kx) .
\end{equation*}
\end{example}

\begin{example}
Temukan solusi $y'' - 6 y' + 13 y = 0$, $y(0) = 0$, $y'(0) = 10$.

Persamaan karakteristiknya adalah $r^2 - 6 r + 13 = 0$. Dengan melengkapkan kuadrat, diperoleh ${(r-3)}^2 + 2^2 = 0$, sehingga akarnya $r = 3 \pm 2i$. Berdasarkan teorema, solusi umumnya
\begin{equation*}
y = C_1 e^{3x} \cos (2x) + C_2 e^{3x} \sin (2x) .
\end{equation*}
Untuk menemukan solusi yang memenuhi syarat awal, mula-mula substitusikan nol untuk memperoleh
\begin{equation*}
0 = y(0) = C_1 e^{0} \cos 0 + C_2 e^{0} \sin 0  = C_1 .
\end{equation*}
Jadi $C_1 = 0$ dan $y = C_2 e^{3x} \sin (2x)$. Kita diferensialkan,
\begin{equation*}
y' = 3C_2 e^{3x} \sin (2x) + 2C_2 e^{3x} \cos (2x) .
\end{equation*}
Kita kembali menerapkan syarat awal dan memperoleh $10 = y'(0) = 2C_2$, atau $C_2 = 5$. Solusi yang dicari adalah
\begin{equation*}
y = 5 e^{3x} \sin (2x) .
\end{equation*}
\end{example}

\subsection{Latihan}

\begin{exercise}
Temukan solusi umum $2y'' + 2y' -4 y = 0$.
\end{exercise}

\begin{exercise}
Temukan solusi umum $y'' + 9y' - 10 y = 0$.
\end{exercise}

\begin{exercise}
Selesaikan $y'' - 8y' + 16 y = 0$ untuk $y(0) = 2$, $y'(0) = 0$.
\end{exercise}

\begin{exercise}
Selesaikan $y'' + 9y' = 0$ untuk $y(0) = 1$, $y'(0) = 1$.
\end{exercise}

\begin{exercise}
Temukan solusi umum $2y'' + 50y = 0$.
\end{exercise}

\begin{exercise}
Temukan solusi umum $y'' + 6 y' + 13 y = 0$.
\end{exercise}

\begin{exercise}
Temukan solusi umum $y'' = 0$ menggunakan metode bagian ini.
\end{exercise}

\begin{exercise}
Metode bagian ini berlaku untuk persamaan yang ordenya bukan dua. Kita akan melihat orde lebih tinggi nanti. Selesaikan persamaan orde pertama $2y' + 3y = 0$ menggunakan metode bagian ini.
\end{exercise}

\begin{exercise}
Mari kita tinjau kembali persamaan Cauchy--Euler\index{persamaan Cauchy--Euler} dari \exercisevref{sol:eulerex}. Sekarang andaikan ${(b-a)}^2-4ac < 0$. Temukan rumus solusi umum $a x^2 y'' + b x y' + c y = 0$. Petunjuk: Perhatikan bahwa $x^r = e^{r \ln x}$.
\end{exercise}

\begin{exercise}
Temukan solusi $y''-(2\alpha) y' + \alpha^2 y=0$, $y(0) = a$, $y'(0)=b$, dengan $\alpha$, $a$, dan $b$ bilangan real.
\end{exercise}

\begin{exercise}
Susun suatu persamaan yang solusi umumnya $y = C_1 e^{-2x} \cos(3x) + C_2 e^{-2x} \sin(3x)$.
\end{exercise}

\setcounter{exercise}{100}

\begin{exercise}
Temukan solusi umum $y''+4y'+2y=0$.
\end{exercise}
\exsol{%
$y = C_1 e^{(-2+\sqrt{2}) x} + C_2 e^{(-2-\sqrt{2}) x}$
}

\begin{exercise}
Temukan solusi umum $y''-6y'+9y=0$.
\end{exercise}
\exsol{%
$y = C_1 e^{3x} + C_2 x e^{3x}$
}

\begin{exercise}
Temukan solusi $2y''+y'+y=0$, $y(0) = 1$, $y'(0)=-2$.
\end{exercise}
\exsol{%
$y = e^{-x/4} \cos\bigl((\nicefrac{\sqrt{7}}{4})x\bigr)
- \sqrt{7} e^{-x/4} \sin\bigl((\nicefrac{\sqrt{7}}{4})x\bigr)$
}

\begin{exercise}
Temukan solusi $2y''+y'-3y=0$, $y(0) = a$, $y'(0)=b$.
\end{exercise}
\exsol{%
$y = \frac{2(a-b)}{5} \, e^{-3x/2}+\frac{3 a+2 b}{5} \, e^x$
}

\begin{exercise}
Temukan solusi $z''(t) = -2z'(t)-2z(t)$, $z(0) = 2$, $z'(0)= -2$.
\end{exercise}
\exsol{%
$z(t) = 2e^{-t} \cos(t)$
}

\begin{exercise}
Temukan solusi $y''-(\alpha+\beta) y' + \alpha \beta y=0$, $y(0) = a$, $y'(0)=b$, dengan $\alpha$, $\beta$, $a$, dan $b$ bilangan real, serta $\alpha \neq \beta$.
\end{exercise}
\exsol{%
$y = \frac{a \beta-b}{\beta-\alpha} e^{\alpha x} +
\frac{b-a \alpha}{\beta-\alpha} e^{\beta x}$
}

\begin{exercise}
Susun suatu persamaan yang solusi umumnya $y = C_1 e^{3x} + C_2 e^{-2x}$.
\end{exercise}
\exsol{%
$y'' -y'-6y=0$
}

%%%%%%%%%%%%%%%%%%%%%%%%%%%%%%%%%%%%%%%%%%%%%%%%%%%%%%%%%%%%%%%%%%%%%%%%%%%%%%

\sectionnewpage
\section{ODE linear orde tinggi} \label{sec:hol}

%mbxINTROSUBSECTION

\sectionnotes{sedikit lebih dari 1 kuliah\EPref{, \S3.2 dan \S3.3 dalam
\cite{EP}}\BDref{,
\S4.1 dan \S4.2 dalam \cite{BD}}}

%After reading this lecture, it may be good to try
%Project III\index{IODE software!Project III} from the
%IODE website: \url{http://www.math.uiuc.edu/iode/}.
%
%\medskip

Kita mempelajari secara singkat persamaan orde tinggi. Persamaan yang muncul dalam penerapan cenderung berorde kedua. Persamaan orde tinggi memang sesekali muncul, tetapi secara umum dunia di sekitar kita \myquote{berorde kedua.}

Hasil dasar tentang ODE linear orde tinggi pada dasarnya sama seperti untuk persamaan orde kedua, dengan 2 diganti $n$. Konsep penting kebebasan linear agak lebih rumit ketika melibatkan lebih dari dua fungsi. Untuk ODE orde tinggi berkoefisien konstan, metode yang dikembangkan juga agak lebih sulit diterapkan, tetapi kita tidak akan berlama-lama pada kerumitan ini. Metode untuk sistem persamaan linear dari \chapterref{sys:chapter} juga dapat digunakan untuk menyelesaikan persamaan orde tinggi berkoefisien konstan.

Perhatikan persamaan linear homogen umum
\begin{equation} \label{hol:eqlinhom}
y^{(n)} + p_{n-1}(x)y^{(n-1)} + \cdots + p_1(x) y' + p_0(x) y = 0 .
\end{equation}

\begin{theorem}[Superposisi]\index{superposisi}
Jika $y_1$, $y_2$, \ldots, $y_n$ adalah solusi persamaan homogen \eqref{hol:eqlinhom}, maka
\begin{equation*}
y(x) = C_1 y_1(x) + C_2 y_2(x) + \cdots + C_n y_n(x)
\end{equation*}
juga menyelesaikan \eqref{hol:eqlinhom} untuk konstanta sebarang $C_1, C_2, \ldots, C_n$.
\end{theorem}

Artinya, suatu \emph{\myindex{kombinasi linear}} dari solusi \eqref{hol:eqlinhom} merupakan solusi \eqref{hol:eqlinhom}. Terdapat pula teorema keberadaan dan ketunggalan untuk persamaan linear, termasuk yang nonhomogen.

\begin{theorem}[Keberadaan dan ketunggalan]\index{keberadaan dan ketunggalan}
Andaikan $p_0, p_1, \ldots, p_{n-1}$, dan $f$ adalah fungsi kontinu pada suatu interval $I$, $a$ adalah suatu bilangan dalam $I$, dan $b_0, b_1, \ldots, b_{n-1}$ adalah konstanta. Maka persamaan
\begin{equation*} %\label{hol:eqlin}
y^{(n)} + p_{n-1}(x)y^{(n-1)} + \cdots + p_1(x) y' + p_0(x) y = f(x)
\end{equation*}
memiliki tepat satu solusi $y(x)$ yang terdefinisi pada interval $I$ yang sama dan memenuhi syarat awal
\begin{equation*}
y(a) = b_0, \quad y'(a) = b_1, \quad \ldots, \quad y^{(n-1)}(a) = b_{n-1} .
\end{equation*}
\end{theorem}

\subsection{Kebebasan linear}

Ketika kita mempunyai dua fungsi $y_1$ dan $y_2$, kita mengatakan keduanya bebas linear jika yang satu bukan kelipatan yang lain. Gagasan yang sama berlaku untuk $n$ fungsi, meskipun dalam kasus ini lebih mudah dinyatakan sebagai berikut. Fungsi $y_1$, $y_2$, \ldots, $y_n$ \emph{\myindex{bebas linear}} jika persamaan
\begin{equation*}
c_1 y_1 + c_2 y_2 + \cdots + c_n y_n = 0
\end{equation*}
hanya mempunyai solusi trivial $c_1 = c_2 = \cdots = c_n = 0$, dengan persamaan harus berlaku untuk semua $x$. Jika persamaan dapat diselesaikan dengan konstanta $c_1,c_2,\ldots,c_n$, misalnya dengan $c_1 \neq 0$, maka $y_1$ dapat dinyatakan sebagai kombinasi linear dari yang lain. Jika fungsi-fungsi tersebut tidak bebas linear, semuanya disebut \emph{\myindex{bergantung linear}}.

\begin{example}
Verifikasi bahwa $e^x, e^{2x}, e^{3x}$ bebas linear.

Mari kita berikan beberapa cara untuk menunjukkan fakta ini. Banyak buku teks (termasuk \cite{EP} dan \cite{F}) memperkenalkan Wronskian, tetapi sulit melihat mengapa metode itu bekerja dan sebenarnya tidak diperlukan di sini.

Perhatikan
\begin{equation*}
c_1 e^x + c_2 e^{2x} + c_3 e^{3x} = 0.
\end{equation*}
Kita gunakan aturan eksponensial dan tulis $z = e^x$. Jadi $z^2 = e^{2x}$ dan $z^3 = e^{3x}$. Maka
\begin{equation*}
c_1 z + c_2 z^2 + c_3 z^3 = 0.
\end{equation*}
Ruas kiri adalah polinomial derajat tiga dalam $z$. Polinomial tersebut identik nol atau mempunyai paling banyak 3 akar. Karena itu, polinomial tersebut identik nol, $c_1 = c_2 = c_3 = 0$, dan fungsi-fungsi tersebut bebas linear.

Mari coba cara lain. Seperti sebelumnya, tulis
\begin{equation*}
c_1 e^x + c_2 e^{2x} + c_3 e^{3x} = 0.
\end{equation*}
Persamaan ini harus berlaku untuk semua $x$. Bagi seluruhnya dengan $e^{3x}$ untuk memperoleh
\begin{equation*}
c_1 e^{-2x} + c_2 e^{-x} + c_3 = 0.
\end{equation*}
Karena persamaan benar untuk semua $x$, biarkan $x \to \infty$. Setelah mengambil limit, terlihat bahwa $c_3 = 0$. Persamaan menjadi
\begin{equation*}
c_1 e^x + c_2 e^{2x} = 0.
\end{equation*}
Ulangi langkah yang sama!

Bagaimana dengan cara lain lagi? Sekali lagi tulis
\begin{equation*}
c_1 e^x + c_2 e^{2x} + c_3 e^{3x} = 0.
\end{equation*}
Kita dapat mengevaluasi persamaan dan turunannya pada nilai $x$ yang berbeda untuk memperoleh persamaan bagi $c_1$, $c_2$, dan $c_3$. Mula-mula bagi dengan $e^{x}$ demi kesederhanaan.
\begin{equation*}
c_1 + c_2 e^{x} + c_3 e^{2x} = 0.
\end{equation*}
Tetapkan $x=0$ untuk memperoleh $c_1 + c_2 + c_3 = 0$. Kini diferensialkan kedua ruas
\begin{equation*}
c_2 e^{x} + 2 c_3 e^{2x} = 0 .
\end{equation*}
Tetapkan $x=0$ untuk memperoleh $c_2 + 2c_3 = 0$. Bagi lagi dengan $e^x$ dan diferensialkan untuk memperoleh $2 c_3 e^{x} = 0$. Jelas bahwa $c_3$ nol. Kemudian $c_2$ harus nol karena $c_2 = -2c_3$, dan $c_1$ harus nol karena $c_1 + c_2 + c_3 = 0$.

Tidak ada satu cara terbaik. Semua metode ini sepenuhnya sah. Yang penting adalah memahami mengapa fungsi-fungsi tersebut bebas linear.
\end{example}

\begin{example}
Di sisi lain, fungsi $e^x$, $e^{-x}$, dan $\cosh x$ bergantung linear. Cukup terapkan definisi kosinus hiperbolik:
\begin{equation*}
\cosh x = \frac{e^x + e^{-x}}{2}
\qquad
\text{atau}
\qquad
2 \cosh x - e^x - e^{-x} = 0.
\end{equation*}
\end{example}

Setelah mempunyai cukup banyak solusi bebas linear, kita memperoleh solusi umum persamaan homogen, sama seperti untuk persamaan orde kedua.

\begin{theorem}
Jika $y_1$, $y_2$, \ldots, $y_n$ adalah solusi bebas linear dari persamaan homogen \eqref{hol:eqlinhom}, maka solusi umum \eqref{hol:eqlinhom} dapat ditulis sebagai
\begin{equation*}
y(x) = C_1 y_1(x) + C_2 y_2(x) + \cdots + C_n y_n(x) .
\end{equation*}
\end{theorem}

\subsection{ODE orde tinggi berkoefisien konstan}

Untuk persamaan homogen linear orde tinggi berkoefisien konstan, prosedurnya persis sama seperti pada orde kedua. Kita hanya perlu menemukan lebih banyak solusi. Jika persamaan berorde $n$, kita perlu menemukan $n$ solusi bebas linear. Hal ini paling baik dilihat melalui contoh.

\begin{example}
Temukan solusi umum
\begin{equation} \label{hol:cceq1}
y''' - 3 y'' - y' + 3y = 0 .
\end{equation}

Coba $y = e^{rx}$. Substitusi memberikan
\begin{equation*}
\underbrace{r^3 e^{rx}}_{y'''} - 3 \underbrace{r^2 e^{rx}}_{y''} -
\underbrace{r e^{rx}}_{y'} + 3 \underbrace{e^{rx}}_{y} = 0 .
\end{equation*}
Bagi seluruhnya dengan $e^{rx}$. Maka
\begin{equation*}
r^3 - 3 r^2 - r + 3 = 0 .
\end{equation*}
Kini kuncinya adalah menemukan akar. Ada rumus untuk akar polinomial derajat 3 dan 4, tetapi sangat rumit. Tidak ada rumus untuk polinomial berderajat lebih tinggi. Ini bukan berarti akarnya tidak ada. Selalu ada $n$ akar untuk polinomial berderajat $n$. Akar dapat berulang\index{akar berulang} dan dapat kompleks. Komputer cukup baik dalam mencari hampiran akar bagi polinomial berukuran wajar.

Langkah awal yang baik adalah menggambar grafik polinomial dan memeriksa letak nolnya. Kita juga dapat langsung mencoba substitusi. Mulailah dengan bilangan $r=-2,-1,0,1,2,\ldots$ dan lihat apakah salah satunya berhasil (bilangan kompleks juga dapat dicoba). Walaupun tidak berhasil, kita mungkin memperoleh petunjuk letak akar. Misalnya, substitusi $r=-2$ menghasilkan $-15$, sedangkan $r=0$ menghasilkan $3$. Ini berarti ada akar antara $r=-2$ dan $r=0$, karena tandanya berubah. Jika menemukan satu akar, misalnya $r_1$, kita tahu $(r-r_1)$ adalah faktor polinomial. Kemudian pembagian panjang polinomial dapat digunakan.

Strategi yang baik ialah memulai dengan $r=0$, $1$, atau $-1$, karena mudah dihitung. Polinomial kita mempunyai dua akar seperti itu, $r_1 = -1$ dan $r_2 = 1$. Seharusnya ada 3 akar, dan akar terakhir cukup mudah ditemukan. Suku konstan dalam polinomial monik\footnote{Kata monik berarti koefisien pangkat tertinggi $r^d$, dalam kasus kita $r^3$, adalah $1$.} seperti ini adalah hasil kali negatif semua akar, karena $r^3 - 3 r^2 - r + 3 = (r-r_1)(r-r_2)(r-r_3)$. Jadi
\begin{equation*}
3 = (-r_1)(-r_2)(-r_3) = (1)(-1)(-r_3) = r_3 .
\end{equation*}
Anda perlu memeriksa bahwa $r_3 = 3$ memang akar. Jadi $e^{-x}$, $e^{x}$, dan $e^{3x}$ adalah solusi \eqref{hol:cceq1}. Ketiganya bebas linear, seperti mudah diperiksa, dan jumlahnya tepat 3, sebanyak yang diperlukan. Jadi solusi umumnya
\begin{equation*}
y = C_1 e^{-x} + C_2 e^{x} + C_3 e^{3x} .
\end{equation*}

Andaikan diberikan syarat awal $y(0) = 1$, $y'(0) = 2$, dan $y''(0) = 3$. Maka
\begin{equation*}
\begin{aligned}
1 = y(0) & = C_1 + C_2 + C_3 , \\
2 = y'(0) & = -C_1 + C_2 + 3C_3 , \\
3 = y''(0) & = C_1 + C_2 + 9C_3 .
\end{aligned}
\end{equation*}
Solusi dapat ditemukan dengan aljabar sekolah menengah, tetapi merepotkan. Cara yang masuk akal untuk menyelesaikan sistem seperti ini adalah menggunakan aljabar matriks; lihat \sectionref{sec:matrix} atau \appendixref{linalg:appendix}. Untuk saat ini, perhatikan bahwa solusinya $C_1 = -\nicefrac{1}{4}$, $C_2 = 1$, dan $C_3 = \nicefrac{1}{4}$. Solusi khusus ODE tersebut adalah
\begin{equation*}
y = \frac{-1}{4}\, e^{-x} + e^x + \frac{1}{4}\, e^{3x} .
\end{equation*}
\end{example}

Berikutnya, andaikan akar-akarnya real tetapi berulang. Misalkan suatu akar $r$ berulang $k$ kali. Sejalan dengan solusi orde kedua, dan untuk alasan yang sama, kita mempunyai solusi
\begin{equation*}
e^{rx}, \quad xe^{rx}, \quad x^2 e^{rx}, \quad \ldots, \quad x^{k-1} e^{rx} .
\end{equation*}
Kita mengambil kombinasi linear dari solusi-solusi ini untuk menemukan solusi umum.

\begin{example}
Selesaikan
\begin{equation*}
y^{(4)} - 3 y''' + 3 y'' - y' =  0 .
\end{equation*}
Persamaan karakteristiknya
\begin{equation*}
r^4 - 3r^3 + 3r^2 -r = 0 .
\end{equation*}
Dengan inspeksi, $r^4 - 3r^3 + 3r^2 -r = r{(r-1)}^3$. Jadi akar-akar beserta \myindex{multiplisitas}-nya adalah $r = 0, 1, 1, 1$. Dengan demikian solusi umumnya
\begin{equation*}
y = \underbrace{(C_1 + C_2 x + C_3 x^2)\, e^x}_{\text{suku yang berasal dari }
r=1} + \underbrace{C_4}_{\text{dari } r=0} .
\end{equation*}
\end{example}

Kasus akar kompleks serupa dengan persamaan orde kedua. Akar kompleks selalu datang berpasangan $r = \alpha \pm i \beta$. Andaikan terdapat dua akar kompleks demikian, masing-masing berulang $k$ kali. Solusi yang bersesuaian adalah
\begin{equation*}
( C_0 + C_1 x + \cdots + C_{k-1} x^{k-1} ) \, e^{\alpha x} \cos (\beta x)
+
( D_0 + D_1 x + \cdots + D_{k-1} x^{k-1} ) \, e^{\alpha x} \sin (\beta x) ,
\avoidbreak
\end{equation*}
di mana $C_0$, \ldots, $C_{k-1}$, $D_0$, \ldots, $D_{k-1}$ adalah konstanta sebarang.

\begin{example}
Selesaikan
\begin{equation*}
y^{(4)} - 4 y''' + 8 y'' - 8 y' + 4y = 0 .
\end{equation*}
Persamaan karakteristiknya
\begin{align*}
r^4 - 4 r^3 + 8 r^2 - 8 r + 4 & = 0 , \\
{(r^2-2r+2)}^2 & = 0 , \\
{\bigl({(r-1)}^2+1\bigr)}^2 & = 0 .
\end{align*}
Jadi akarnya $1 \pm i$, keduanya bermultiplisitas 2. Maka solusi umum ODE tersebut
\begin{equation*}
y =
( C_1 + C_2 x ) \, e^{x} \cos x
+
( C_3 + C_4 x ) \, e^{x} \sin x .
\end{equation*}
Cara kita menyelesaikan persamaan karakteristik di atas sebenarnya dengan menebak atau inspeksi. Secara umum tidak semudah itu. Kita juga dapat meminta komputer atau kalkulator canggih untuk mencari akarnya.
\end{example}

%FIXME: the operator stuff?

\subsection{Latihan}

\begin{exercise}
Temukan solusi umum $y''' - y'' + y' - y = 0$.
\end{exercise}

\begin{exercise}
Temukan solusi umum $y^{(4)} - 5 y''' + 6 y'' = 0$.
\end{exercise}

\begin{exercise}
Temukan solusi umum $y''' + 2 y'' + 2 y' = 0$.
\end{exercise}

\begin{exercise}
Andaikan persamaan karakteristik suatu ODE adalah ${(r-1)}^2{(r-2)}^2 = 0$.
\begin{tasks}
\task Temukan persamaan diferensial semacam itu.
\task Temukan solusi umumnya.
\end{tasks}
\end{exercise}

\begin{exercise} \label{hol:eqfromsolex}
Andaikan suatu persamaan orde keempat mempunyai solusi $y = 2 e^{4x} x \cos x$.
\begin{tasks}
\task Temukan persamaan semacam itu.
\task Temukan syarat awal yang dipenuhi solusi tersebut.
\end{tasks}
\end{exercise}

\begin{exercise}
Temukan solusi umum persamaan dalam \exerciseref{hol:eqfromsolex}.
\end{exercise}

\begin{exercise}
Misalkan $f(x) = e^x - \cos x$, $g(x) = e^x + \cos x$, dan $h(x) = \cos x$. Apakah $f(x)$, $g(x)$, dan $h(x)$ bebas linear? Jika ya, tunjukkan; jika tidak, temukan kombinasi linear yang sesuai.
\end{exercise}

\begin{exercise}
Misalkan $f(x) = 0$, $g(x) = \cos x$, dan $h(x) = \sin x$. Apakah $f(x)$, $g(x)$, dan $h(x)$ bebas linear? Jika ya, tunjukkan; jika tidak, temukan kombinasi linear yang sesuai.
\end{exercise}

\begin{exercise}
Apakah $x$, $x^2$, dan $x^4$ bebas linear? Jika ya, tunjukkan; jika tidak, temukan kombinasi linear yang menghasilkan $0$.
\end{exercise}

\begin{exercise}
Apakah $e^x$, $xe^x$, dan $x^2e^x$ bebas linear? Jika ya, tunjukkan; jika tidak, temukan kombinasi linear yang menghasilkan $0$.
\end{exercise}

\begin{exercise}
Temukan suatu persamaan yang mempunyai $y=xe^{-2x}\sin(3x)$ sebagai solusi.
\end{exercise}

\setcounter{exercise}{100}

\begin{exercise}
Temukan solusi umum $y^{(5)}-y^{(4)}=0$.
\end{exercise}
\exsol{%
$y=C_1 e^x +C_2 x^3 + C_3 x^2 +C_4 x + C_5$
}

\begin{exercise}
\pagebreak[2]
Andaikan persamaan karakteristik suatu persamaan diferensial orde ketiga mempunyai akar $\pm 2i$ dan 3.
\begin{tasks}
\task Apa persamaan karakteristiknya?
\task Temukan persamaan diferensial yang bersesuaian.
\task Temukan solusi umumnya.
\end{tasks}
\end{exercise}
\exsol{%
a) $r^3-3r^2+4r-12 = 0$
\quad
b) $y'''-3y''+4y'-12y = 0$
\quad
c) $y = C_1 e^{3x} + C_2 \sin(2x) + C_3 \cos(2x)$
}

\begin{exercise}
Selesaikan $1001y'''+3.2y''+\pi y'-\sqrt{4} y = 0$, $y(0)=0$, $y'(0) = 0$, $y''(0) = 0$.
\end{exercise}
\exsol{%
$y=0$
}

\begin{exercise}
Apakah $e^{x}$, $e^{x+1}$, $e^{2x}$, $\sin(x)$ bebas linear? Jika ya, tunjukkan; jika tidak, temukan kombinasi linear yang menghasilkan $0$.
\end{exercise}
\exsol{%
Tidak. $e^1 e^x -  e^{x+1} = 0$.
}

\begin{exercise}
Apakah $\sin(x)$, $x$, $x\sin(x)$ bebas linear? Jika ya, tunjukkan; jika tidak, temukan kombinasi linear yang menghasilkan $0$.
\end{exercise}
\exsol{%
Ya. (Petunjuk: Pertama, perhatikan bahwa $\sin(x)$ terbatas. Kemudian perhatikan bahwa $x$ dan $x\sin(x)$ tidak mungkin saling merupakan kelipatan.)
}

\begin{exercise}
Temukan suatu persamaan yang mempunyai $y=\cos(x)$, $y=\sin(x)$, $y=e^x$ sebagai solusi.
\end{exercise}
\exsol{%
$y'''-y''+y'-y=0$
}

%%%%%%%%%%%%%%%%%%%%%%%%%%%%%%%%%%%%%%%%%%%%%%%%%%%%%%%%%%%%%%%%%%%%%%%%%%%%%%

\sectionnewpage
\section{Getaran mekanis} \label{sec:mv}

%mbxINTROSUBSECTION

\sectionnotes{2 kuliah\EPref{, \S3.4 dalam \cite{EP}}\BDref{,
\S3.7 dalam \cite{BD}}}

Mari kita lihat beberapa penerapan persamaan orde kedua berkoefisien konstan.

\subsection{Beberapa contoh}

\begin{mywrapfigsimp}{2.0in}{2.3in}
\noindent
\inputpdft{massfigforce}{%
%note that this diagram appears more than once in the book
Diagram yang menunjukkan massa berlabel m, terhubung di sebelah kiri ke pegas berlabel k yang terhubung ke dinding di sebelah kiri. Panah gaya ke kanan pada massa berlabel F dari t. Panah di bawah massa yang menunjukkan perpindahan mengarah ke kanan dan berlabel x. `Lantai' diagram berlabel `redaman c'.}
\end{mywrapfigsimp}
Contoh pertama kita adalah massa pada pegas. Andaikan massa $m > 0$ (dalam kilogram) dihubungkan ke dinding tetap oleh pegas dengan konstanta pegas $k > 0$ (dalam newton per meter). Mungkin terdapat gaya eksternal $F(t)$ (dalam newton) yang bekerja pada massa. Di sini $t$ adalah waktu dalam detik. Terakhir, terdapat gesekan yang diukur oleh $c \geq 0$ (dalam newton-detik per meter) saat massa meluncur di lantai (atau mungkin terhubung ke peredam).

Misalkan $x$ adalah perpindahan massa dalam meter ($x=0$ adalah posisi diam), dengan $x$ membesar ke kanan (menjauhi dinding). Berdasarkan \myindex{hukum Hooke}, gaya yang diberikan pegas sebanding dengan kompresi pegas. Jadi gaya itu adalah $kx$ dalam arah negatif. Demikian pula, gaya gesek sebanding dengan kecepatan massa. \myindex{Hukum kedua Newton} mengatakan bahwa gaya sama dengan massa kali percepatan. Jadi $mx'' = F(t)-cx'-kx$ atau
\begin{equation*}
mx'' + cx' + kx = F(t) .
\end{equation*}
Persamaan ini adalah ODE linear orde kedua berkoefisien konstan\@. Kita mengatakan geraknya
\begin{enumerate}[(i)]
\item \emph{dipaksa\index{gerak dipaksa}} jika $F \not\equiv 0$ (jika $F$ tidak identik nol),
\item \emph{tak dipaksa\index{gerak tak dipaksa}} atau \emph{bebas\index{gerak bebas}} jika $F \equiv 0$ (jika $F$ identik nol),
\item \emph{teredam\index{gerak teredam}} jika $c > 0$, dan
\item \emph{tak teredam\index{gerak tak teredam}} jika $c = 0$.
\end{enumerate}

Persamaan ini muncul dalam banyak penerapan walaupun awalnya mungkin tidak tampak demikian. Banyak situasi nyata dapat disederhanakan menjadi massa pada pegas. Misalnya, lompat bungee pada dasarnya adalah sistem massa dan pegas (Anda adalah massanya). Sebaiknya seseorang menghitung matematikanya sebelum Anda melompat dari jembatan, bukan? Berikut dua contoh lain.

\medskip

%5 is the number of lines, must be adjusted
\begin{mywrapfigsimp}[5]{1.35in}{1.65in}
\noindent
\inputpdft{mv-rlc}{%
Diagram rangkaian. Dengan mengikuti diagram berlawanan arah jarum jam, kita mulai di kiri dengan sumber listrik berlabel E, lalu resistor di bawah berlabel R, induktor di kanan berlabel L, kapasitor di atas berlabel C, kemudian kembali terhubung ke sumber listrik.}
\end{mywrapfigsimp}
Berikut contoh untuk insinyur listrik. Perhatikan \myindex{rangkaian RLC} pada gambar. Terdapat resistor dengan resistansi $R$ ohm, induktor dengan induktansi $L$ henry, dan kapasitor dengan kapasitansi $C$ farad. Ada pula sumber listrik (misalnya baterai) dengan tegangan $E(t)$ volt pada waktu $t$ (diukur dalam detik). Misalkan $Q(t)$ adalah muatan (dalam coulomb) pada kapasitor dan $I(t)$ adalah arus (dalam ampere) dalam rangkaian. Hubungan keduanya adalah $Q' = I$. Berdasarkan prinsip dasar, diperoleh $L I' + RI + \nicefrac{Q}{C} = E$. Kita diferensialkan untuk memperoleh
\begin{equation*}
L I''(t) + R I'(t) + \frac{1}{C} I(t) = E'(t) .
\end{equation*}
Ini adalah persamaan linear nonhomogen orde kedua berkoefisien konstan. Karena $L, R$, dan $C$ semuanya positif, rangkaian ini berperilaku seperti sistem massa dan pegas. Posisi massa digantikan oleh arus; massa oleh induktansi; redaman oleh resistansi; dan konstanta pegas oleh kebalikan kapasitansi. Perubahan tegangan menjadi fungsi pemaksa---untuk tegangan konstan, geraknya tak dipaksa.

\medskip

%9 is the number of lines, must be adjusted
\begin{mywrapfigsimp}[9]{1.8in}{2.16in}
\noindent
\inputpdft{mv-pend-deriv}{%
Diagram bandul. Bandul tampak terikat pada titik di atas dan bergerak sepanjang busur lingkaran. Bagian bawah bandul berlabel m, dan sudut antara bandul dan garis vertikal berlabel theta. Panjang bandul berlabel L. Terdapat panah gaya pada m sepanjang busur menjauhi pusat berlabel m kali L kali turunan kedua theta, panah gaya berlawanan arah (menuju pusat) berlabel m kali g kali sin theta, dan panah vertikal ke bawah dari massa berlabel m kali g. Dua panah terakhir membentuk segitiga siku-siku dengan panah ke bawah sebagai hipotenusa.}
\end{mywrapfigsimp}
Contoh berikutnya hanya mendekati perilaku sistem massa dan pegas. Andaikan massa $m$ kilogram tergantung pada bandul sepanjang $L$ meter. Kita mencari persamaan untuk sudut $\theta(t)$ (dalam radian). Misalkan $g$ adalah percepatan gravitasi dalam meter per detik kuadrat. Mari turunkan persamaan $\theta$ menggunakan \myindex{hukum kedua Newton}: Gaya sama dengan massa kali percepatan. Percepatannya $L \theta''$ dan massanya $m$. Jadi $mL\theta''$ harus sama dengan komponen tangensial gaya gravitasi, yaitu $m g \sin \theta$ dalam arah berlawanan. Jadi $mL\theta'' = -mg \sin \theta$. Secara menarik, $m$ saling menghapus dan kita memperoleh
\begin{equation*}
\theta'' + \frac{g}{L} \sin \theta = 0 .
\end{equation*}

Kini kita membuat hampiran. Untuk $\theta$ kecil, secara hampiran $\sin \theta \approx \theta$. Hal ini terlihat dari grafik. Dalam \figurevref{mv:sinthetafig}, untuk kira-kira $-0.5 < \theta < 0.5$ (dalam radian), grafik $\sin \theta$ dan $\theta$ hampir sama.

\begin{mywrapfig}{3.25in}
\capstart
\diffyincludegraphics{width=3in}{width=4.5in}{mv-sintheta}{%
Grafik sin theta dan theta untuk theta antara minus 1 dan 1. Grafik theta adalah garis lurus dari (minus 1, minus 1) ke (1,1), sedangkan grafik sin theta hampir sama, tetapi berawal sedikit di atas garis lurus di kiri dan berakhir sedikit di bawahnya di kanan.}
\caption{Grafik $\sin \theta$ dan $\theta$ (dalam radian).\label{mv:sinthetafig}}
\end{mywrapfig}

Jadi ketika ayunan kecil, $\theta$ kecil, dan perilakunya dapat dimodelkan dengan persamaan linear yang lebih sederhana
\begin{equation*}
\theta'' + \frac{g}{L} \theta = 0 .
\end{equation*}
Galat hampiran ini terakumulasi. Setelah waktu yang lama, keadaan sistem nyata dapat berbeda secara berarti dari solusi kita. Selain itu, akan kita lihat bahwa dalam sistem massa-pegas amplitudo tidak bergantung pada periode. Hal ini tidak benar untuk bandul. Meskipun demikian, untuk rentang waktu yang cukup pendek dan ayunan kecil (hanya sudut $\theta$ kecil), hampirannya cukup baik.

Dalam masalah nyata, penyederhanaan semacam ini sering diperlukan. Kita harus memahami matematika dan fisika situasinya untuk menilai apakah penyederhanaan sah dalam konteks pertanyaan yang hendak dijawab.

\subsection{Gerak bebas tak teredam}

Dalam bagian ini kita hanya meninjau gerak bebas atau tak dipaksa, karena kita belum tahu cara menyelesaikan persamaan nonhomogen. Kita mulai dengan gerak \myindex{tak teredam}, ketika $c=0$. Persamaannya
\begin{equation*}
mx'' + kx = 0 .
\end{equation*}
Bagi dengan $m$ dan tetapkan $\omega_0 = \sqrt{\nicefrac{k}{m}}$ untuk menulis ulang persamaan sebagai
\begin{equation*}
x'' + \omega_0^2 x = 0 .
\end{equation*}
Solusi umumnya
\begin{equation*}
x(t) = A \cos (\omega_0 t) + B \sin (\omega_0 t) .
\end{equation*}
Berdasarkan identitas trigonometri
\begin{equation*}
A \cos (\omega_0 t) + B \sin (\omega_0 t) =
C \cos ( \omega_0 t - \gamma ) ,
\end{equation*}
untuk dua konstanta lain $C$ dan $\gamma$. Diperoleh $A = C \cos \gamma$ dan $B = C \sin \gamma$, sehingga mudah dihitung bahwa $C= \sqrt{A^2 + B^2}$ dan $\tan \gamma = \nicefrac{B}{A}$. Jadi kita ambil $C$ dan $\gamma$ sebagai konstanta sebarang dan tulis $x(t) = C \cos ( \omega_0 t - \gamma )$.

\begin{exercise}
Benarkan identitas $A \cos (\omega_0 t) + B \sin (\omega_0 t) = C \cos ( \omega_0 t - \gamma )$ dan verifikasi persamaan $C = \sqrt{A^2+B^2}$ serta $\tan \gamma = \frac{B}{A}$. Petunjuk: Mulailah dengan $\cos (\alpha-\beta) = \cos (\alpha) \cos (\beta) + \sin (\alpha)\sin (\beta)$ dan kalikan dengan $C$. Lalu apa yang semestinya menjadi $\alpha$ dan $\beta$?
\end{exercise}

Walaupun bentuk pertama dengan $A$ dan $B$ lebih mudah digunakan untuk menentukan syarat awal, bentuk kedua lebih alami. Konstanta $C$ dan $\gamma$ mempunyai interpretasi fisik yang baik. Tuliskan solusi sebagai
\begin{equation*}
x(t) = C \cos ( \omega_0 t - \gamma ) .
\end{equation*}
Ini adalah osilasi berfrekuensi tunggal (gelombang sinus). \emph{\myindex{Amplitudo}}-nya adalah $C$, $\omega_0$ adalah \emph{\myindex{frekuensi}} (sudut)\index{frekuensi sudut}, dan $\gamma$ adalah apa yang disebut \emph{\myindex{pergeseran fase}}. Pergeseran fase hanya menggeser grafik ke kiri atau kanan. Kita menyebut $\omega_0$ sebagai \emph{\myindex{frekuensi (sudut) alami}}. Keseluruhan susunan ini disebut \emph{\myindex{gerak harmonik sederhana}}.

Mari berhenti sejenak untuk menjelaskan kata \emph{sudut} sebelum kata \emph{frekuensi}. Satuan $\omega_0$ adalah radian per satuan waktu, bukan siklus per satuan waktu seperti ukuran frekuensi biasa. Karena satu siklus adalah $2\pi$ radian, frekuensi biasa diberikan oleh $\frac{\omega_0}{2\pi}$. Ini sekadar perkara tempat kita menaruh konstanta $2\pi$, dan itu soal selera.

\emph{\myindex{Periode}} gerak adalah satu dibagi frekuensi (dalam siklus per satuan waktu), sehingga bernilai $\frac{2\pi}{\omega_0}$. Itulah waktu untuk menyelesaikan satu siklus penuh.

\begin{example}
Andaikan $m=\unit[2]{kg}$ dan $k=\unitfrac[8]{N}{m}$. Seluruh susunan massa dan pegas berada di atas truk yang melaju dengan \unitfrac[1]{m}{s}. Truk menabrak lalu berhenti. Massa ditahan 0.5 meter di depan posisi diam. Saat tabrakan, massa terlepas. Artinya, massa kini bergerak maju dengan \unitfrac[1]{m}{s}, sementara ujung pegas lainnya tetap. Jadi massa mulai berosilasi. Berapa frekuensi osilasinya? Berapa amplitudonya? Satuan yang digunakan adalah satuan mks\index{satuan mks} (meter-kilogram-detik).

Susunan ini berarti pada saat tabrakan massa berada setengah meter dalam arah positif dan, relatif terhadap dinding tempat pegas dipasang, bergerak maju (arah positif) dengan \unitfrac[1]{m}{s}. Ini memberikan syarat awal. Jadi persamaan beserta syarat awal adalah
\begin{equation*}
2 x'' + 8 x = 0 , \qquad x(0) = 0.5, \qquad x'(0) = 1.
\end{equation*}
Kita hitung $\omega_0 = \sqrt{\nicefrac{k}{m}} = \sqrt{4} = 2$. Jadi frekuensi sudutnya 2. Frekuensi biasa dalam Hertz (siklus per detik) adalah $\nicefrac{2}{2\pi} = \nicefrac{1}{\pi} \approx 0.318$.

Solusi umumnya
\begin{equation*}
x(t) = A \cos (2t) + B \sin (2t) .
\end{equation*}
Syarat $x(0) = 0.5$ berarti $A = 0.5$. Lalu $x'(t) = - 2(0.5) \sin (2t) + 2B \cos (2t)$. Dari $x'(0) = 1$, diperoleh $B = 0.5$. Jadi amplitudonya $C = \sqrt{A^2+B^2} = \sqrt{0.25+0.25} = \sqrt{0.5} \approx 0.707$. Solusinya
\begin{equation*}
x(t) = 0.5 \cos (2t) + 0.5 \sin (2t) .
\end{equation*}
Grafik $x(t)$ ditampilkan dalam \figurevref{mv:undampedfig}.
\end{example}

%15 is the number of lines, must be adjusted
\begin{mywrapfig}[15]{3.25in}
\capstart
\diffyincludegraphics{width=3in}{width=4.5in}{mv-undamped}{%
Grafik fungsi untuk t dari 0 sampai 10, berupa kurva sinusoidal antara kira-kira minus 0.707 dan 0.707, dimulai di (0,0.5) dan naik, serta membentuk sedikit lebih dari 3 siklus dalam gambar.}
\caption{Osilasi sederhana tak teredam.\label{mv:undampedfig}}
\end{mywrapfig}

Secara umum, untuk gerak bebas tak teredam, solusi berbentuk
\begin{equation*}
x(t) = A \cos (\omega_0 t) + B \sin (\omega_0 t) ,
\end{equation*}
bersesuaian dengan syarat awal $x(0) = A$ dan $x'(0) = \omega_0 B$. Mudah menemukan $A$ dan $B$ dari syarat awal. Amplitudo dan pergeseran fase dihitung dari $A$ dan $B$. Dalam contoh, kita menemukan amplitudo $C$. Bagaimana dengan pergeseran fase $\gamma$? Kita tahu $\tan \gamma = \nicefrac{B}{A} = 1$. Arkustangen $1$ adalah $\nicefrac{\pi}{4}$ atau sekitar $0.785$. Kita harus memeriksa apakah $\nicefrac{\pi}{4}$ memberikan kuadran yang benar---apakah sudut dari sumbu $x$ positif berada dalam kuadran yang sama dengan titik $(A,B)$---jika tidak, tambahkan $\pi$. Alasannya $A = C \cos \gamma$ dan $B = C \sin \gamma$. Karena $A$ dan $B$ positif, $\gamma$ harus di kuadran pertama. Karena $\nicefrac{\pi}{4}$ radian berada di kuadran pertama, $\gamma = \nicefrac{\pi}{4}$.

Catatan: Banyak kalkulator dan perangkat lunak komputer tidak hanya memiliki fungsi \texttt{atan}\index{atan} untuk arkustangen, tetapi juga fungsi yang kadang disebut \texttt{atan2}\index{atan2}. Fungsi ini menerima dua argumen, $B$ dan $A$, lalu mengembalikan $\gamma$ dalam kuadran yang benar.

\subsection{Gerak bebas teredam}

%mbxINTROSUBSUBSECTION

Sekarang kita berfokus pada gerak \myindex{teredam}. Tulis ulang persamaan
\begin{equation*}
m x'' + c x' + kx = 0,
\end{equation*}
sebagai
\begin{equation*}
x'' + 2p x' + \omega_0^2 x = 0,
\end{equation*}
di mana
\begin{equation*}
\omega_0 = \sqrt{\frac{k}{m}}, \qquad p = \frac{c}{2m} .
\end{equation*}
Persamaan karakteristiknya
\begin{equation*}
r^2 + 2 pr + \omega_0^2 = 0 .
\end{equation*}
Berdasarkan rumus kuadrat, akarnya
\begin{equation*}
r = -p \pm \sqrt{p^2 - \omega_0^2} .
\end{equation*}
Bentuk solusi bergantung pada apakah akarnya kompleks atau real. Akar real diperoleh jika dan hanya jika bilangan berikut taknegatif:
\begin{equation*}
p^2 - \omega_0^2 = {\left( \frac{c}{2m} \right)}^2 - \frac{k}{m}
= \frac{c^2 - 4km}{4m^2} .
\end{equation*}
Tanda $p^2-\omega_0^2$ sama dengan tanda $c^2 - 4km$. Jadi akar real diperoleh jika dan hanya jika $c^2-4km$ taknegatif, atau $c^2 \geq 4km$.

\subsubsection{Redaman berlebih}

%15 is the number of lines, must be adjusted
%mbxSTARTIGNORE
\begin{mywrapfig}[15]{3.25in}
\capstart
\diffyincludegraphics{width=3in}{width=4.5in}{mv-overdamped}{%
Grafik 3 fungsi untuk t antara 0 dan 100. Ketiganya bermula di (0,1) dan mendekati sumbu horizontal hingga hampir tidak dapat dibedakan darinya di sisi kanan gambar. Fungsi pertama mula-mula turun, memotong sumbu, lalu naik kembali dan mendekatinya; fungsi kedua bermula mendatar, kemudian berbelok turun dan mendekati sumbu; fungsi ketiga mula-mula naik, lalu cepat berbelok turun dan mendekati sumbu.}
\caption{Gerak teredam berlebih untuk beberapa syarat awal berbeda.\label{mv:overdampedfig}}
\end{mywrapfig}
%mbxENDIGNORE
%
% make sure the MBX below is synced!
%

Ketika $c^2 - 4km > 0$, sistem \emph{\myindex{teredam berlebih}}. Dalam kasus ini terdapat dua akar real berbeda $r_1$ dan $r_2$. Karena $\sqrt{p^2 - \omega_0^2}$ selalu kurang dari $p$, ekspresi akar $-p \pm \sqrt{p^2 - \omega_0^2}$ negatif untuk kedua tanda. Jadi kedua akar negatif.

Solusinya
\begin{equation*}
x(t) = C_1 e^{r_1 t} + C_2 e^{r_2 t} .
\end{equation*}
Karena $r_1, r_2$ negatif, $x(t) \to 0$ saat $t \to \infty$. Jadi massa cenderung menuju posisi diam ketika waktu menuju tak hingga. Untuk beberapa contoh grafik dengan syarat awal berbeda, lihat \figurevref{mv:overdampedfig}.

%mbxlatex \begin{myfig}
%mbxlatex \diffyincludegraphics{width=3in}{width=4.5in}{mv-overdamped}{%
%mbxlatex A plot of 3 functions for t between 0 and 100.  All three graphs start at
%mbxlatex (0,1) and they all approach the horizontal axis being almost
%mbxlatex indistinguishable from it by the right side of the picture.  The first one
%mbxlatex starts downwards, crosses the axis and then goes back up and approaches it,
%mbxlatex the second starts out horizontally but turns downwards and approaches the
%mbxlatex axis, and the third starts upwards, then turns downwards quickly and
%mbxlatex approaches the axis.}
%mbxlatex \caption{Overdamped motion for several different initial conditions.\label{mv:overdampedfig}}
%mbxlatex \end{myfig}

Tidak terjadi osilasi. Bahkan, grafik memotong sumbu $t$ paling banyak sekali. Untuk melihat alasannya, coba selesaikan $0 = C_1 e^{r_1 t} + C_2 e^{r_2 t}$. Maka $C_1 e^{r_1 t} = - C_2 e^{r_2 t}$ dan aturan eksponen memberikan
\begin{equation*}
\frac{-C_1}{C_2} = e^{(r_2-r_1) t} .
\end{equation*}
Persamaan ini mempunyai paling banyak satu solusi $t \geq 0$; grafik $e^{(r_2-r_1) t}$ selalu naik atau selalu turun. Untuk beberapa syarat awal, grafik tidak pernah memotong sumbu $t$, seperti tampak pada contoh grafik.

\begin{example}
Andaikan massa dilepaskan dari keadaan diam, yaitu $x(0) = x_0$ dan $x'(0) = 0$. Maka
\begin{equation*}
x(t) = \frac{x_0}{r_1-r_2} \left(r_1 e^{r_2 t} - r_2 e^{r_1 t} \right) .
\end{equation*}
Mudah dilihat bahwa ini memenuhi syarat awal.
\end{example}

\subsubsection{Redaman kritis}

Ketika $c^2 - 4km = 0$, sistem \emph{\myindex{teredam kritis}}. Dalam kasus ini terdapat satu akar bermultiplisitas 2, yaitu $-p$. Solusinya
\begin{equation*}
x(t) = C_1 e^{-pt} + C_2 t e^{-pt} .
\end{equation*}
Perilaku sistem teredam kritis sangat mirip dengan sistem teredam berlebih. Bagaimanapun, dalam arti tertentu sistem teredam kritis adalah limit dari sistem teredam berlebih. Karena persamaan ini sebenarnya hanya hampiran dunia nyata, dalam kenyataan kita tidak pernah tepat teredam kritis; keadaan itu hanya dapat dicapai secara teoretis. Kita selalu sedikit teredam kurang atau sedikit teredam berlebih. Tidak perlu terlalu berlama-lama pada redaman kritis.

\subsubsection{Redaman kurang}

%13 is the number of lines, must be adjusted
%mbxSTARTIGNORE
\begin{mywrapfig}[13]{3.25in}
\capstart
\diffyincludegraphics{width=3in}{width=4.5in}{mv-underdamped}{%
Grafik untuk t antara 0 dan 30. Ditampilkan dua kurva eksponensial meluruh, satu bermula di (0,1) dan satu di (0, minus 1). Kurva mirip sinusoidal berkelok di antara keduanya, bermula pada kurva atas, turun ke kurva bawah, kembali naik, dan seterusnya.}
\caption{Gerak teredam kurang dengan kurva selubung ditampilkan.\label{mv:underdampedfig}}
\end{mywrapfig}
%mbxENDIGNORE
%
% make sure the MBX below is synced!
%
Ketika $c^2 - 4km < 0$, sistem \emph{\myindex{teredam kurang}}. Dalam kasus ini, akarnya kompleks.
\begin{equation*}
\begin{split}
r & = -p \pm \sqrt{p^2 - \omega_0^2} \\
& = -p \pm \sqrt{-1}\sqrt{\omega_0^2 - p^2} \\
& = -p \pm i \omega_1 ,
\end{split}
\end{equation*}
di mana $\omega_1 =\sqrt{\omega_0^2 - p^2}$. Solusinya
\begin{equation*}
x(t) = e^{-pt} \bigl( A \cos (\omega_1 t) + B \sin (\omega_1 t) \bigr) ,
\end{equation*}
atau
\begin{equation*}
x(t) = C e^{-pt} \cos ( \omega_1 t - \gamma ) .
\end{equation*}
Contoh grafik diberikan dalam \figurevref{mv:underdampedfig}. Perhatikan bahwa tetap berlaku $x(t) \to 0$ saat $t \to \infty$.

%mbxlatex \begin{myfig}
%mbxlatex \diffyincludegraphics{width=3in}{width=4.5in}{mv-underdamped}{%
%mbxlatex A plot for t between 0 and 30.  Two decaying exponential curves one starting it (0,1)
%mbxlatex and one starting at (0, minus 1) are shown.  A sinusoidal-like curve winds between
%mbxlatex them, starting at the top curve, going down the bottom curve, then back up
%mbxlatex and so on.}
%mbxlatex \caption{Underdamped motion with the envelope curves shown.\label{mv:underdampedfig}}
%mbxlatex \end{myfig}

Gambar juga menunjukkan \emph{\myindex{kurva selubung}} $C e^{-pt}$ dan $-C e^{-pt}$. Solusi adalah garis berosilasi di antara kedua kurva selubung. Kurva selubung memberikan amplitudo maksimum osilasi pada setiap saat. Misalnya, jika Anda melompat bungee, Anda benar-benar berkepentingan menghitung kurva selubung agar kepala tidak membentur beton.
%
Pergeseran fase $\gamma$ menggeser osilasi ke kiri atau kanan, tetapi tetap di dalam kurva selubung (kurva selubung tidak berubah jika $\gamma$ berubah).

Perhatikan bahwa \emph{\myindex{frekuensi semu}} sudut\footnote{Kita tidak menyebut $\omega_1$ sebagai frekuensi karena solusinya bukan benar-benar fungsi periodik.} $\omega_1$ menjadi lebih kecil ketika redaman $c$ (dan karenanya $p$) membesar. Ini masuk akal. Pertama, jika redaman hanya sedikit diubah, kita tidak mengharapkan perilaku solusi berubah drastis. Kedua, jika $c$ terus diperbesar, pada titik tertentu solusi harus mulai menyerupai solusi redaman kritis atau redaman berlebih, tanpa osilasi. Saat $c$ membesar dan $c^2$ mendekati $4km$, $\omega_1$ mendekati $0$.
%
Di sisi lain, ketika $c$ mengecil, $\omega_1$ mendekati $\omega_0$ ($\omega_1$ selalu lebih kecil dari $\omega_0$), dan solusi semakin menyerupai gerak periodik mantap pada kasus tak teredam. Kurva selubung semakin datar saat $c$ (dan $p$) menuju $0$.

\subsection{Latihan}

\begin{samepage}
\begin{exercise} \label{mv:ex1}
Perhatikan sistem massa dan pegas dengan massa $m=2$, konstanta pegas $k=3$, dan konstanta redaman $c=1$.
\begin{tasks}
\task Susun dan temukan solusi umum sistem.
\task Apakah sistem teredam kurang, teredam berlebih, atau teredam kritis?
\task Jika sistem tidak teredam kritis, temukan $c$ yang membuatnya teredam kritis.
\end{tasks}
\end{exercise}
\end{samepage}

\begin{exercise}
Kerjakan \exerciseref{mv:ex1} untuk $m=3$, $k=12$, dan $c=12$.
\end{exercise}

\begin{exercise} \label{mv:exwt1}
Dengan satuan mks (meter-kilogram-detik)\index{satuan mks}, andaikan Anda mempunyai pegas dengan konstanta \unitfrac[4]{N}{m}. Anda ingin menggunakannya untuk menimbang benda. Andaikan tanpa gesekan. Letakkan massa pada pegas dan gerakkan.
\begin{tasks}
\task Anda menghitung dan menemukan frekuensi \unit[0.8]{Hz} (siklus per detik). Berapa massanya?
\task Temukan rumus massa $m$ jika diberikan frekuensi $\omega$ dalam \unit{Hz}.
\end{tasks}
\end{exercise}

\begin{exercise}
Andaikan kita menambahkan kemungkinan gesekan pada \exerciseref{mv:exwt1}. Andaikan pula konstanta pegas tidak diketahui, tetapi tersedia dua beban acuan \unit[1]{kg} dan \unit[2]{kg} untuk mengalibrasi susunan. Gerakkan masing-masing pada pegas dan ukur frekuensinya. Untuk beban \unit[1]{kg}, terukur \unit[1.1]{Hz}; untuk beban \unit[2]{kg}, terukur \unit[0.8]{Hz}.
\begin{tasks}
\task Temukan $k$ (konstanta pegas) dan $c$ (konstanta redaman).
\task Temukan rumus massa dalam frekuensi Hz. \emph{Perhatikan bahwa mungkin terdapat lebih dari satu massa untuk frekuensi tertentu.}
\task Untuk benda tak dikenal terukur \unit[0.2]{Hz}; berapa massanya? Andaikan diketahui bahwa massa benda lebih dari satu kilogram.
\end{tasks}
\end{exercise}

\begin{exercise}
Andaikan Anda ingin mengukur gesekan yang dialami massa \unit[0.1]{kg} ketika meluncur di lantai (Anda ingin menemukan $c$). Tersedia pegas dengan konstanta $k=\unitfrac[5]{N}{m}$. Hubungkan pegas ke massa dan pasang pada dinding. Tarik pegas lalu lepaskan massa. Ditemukan bahwa massa berosilasi dengan frekuensi \unit[1]{Hz}. Berapa gesekannya?
\end{exercise}

\setcounter{exercise}{100}

\begin{exercise}
Massa $2$ kilogram berada pada pegas dengan konstanta pegas $k$ newton per meter tanpa redaman. Andaikan sistem diam dan pada waktu $t=0$ massa ditendang sehingga mulai bergerak 2 meter per detik. Seberapa besar $k$ agar massa tidak bergerak lebih dari 3 meter dari posisi diam?
\end{exercise}
\exsol{%
$k=\unitfrac[\nicefrac{8}{9}]{N}{m}$ (atau lebih besar)
}

\begin{exercise}
Andaikan terdapat rangkaian RLC dengan resistor 100 miliohm (0.1 ohm), induktor berinduktansi 50 milihenry (0.05 henry), kapasitor 5 farad, dan tegangan konstan.
\begin{tasks}
\task Susun persamaan ODE untuk arus $I$.
\task Temukan solusi umum.
\task Selesaikan untuk $I(0) = 10$ dan $I'(0) = 0$.
\end{tasks}
\end{exercise}
\exsol{%
a) $0.05 I'' + 0.1 I' + (\nicefrac{1}{5}) I = 0$
\quad
b) $I = C e^{-t} \cos(\sqrt{3} \, t - \gamma)$
\quad
c) $I = 10 e^{-t} \cos(\sqrt{3} \, t) + \frac{10}{\sqrt{3}} e^{-t}
\sin(\sqrt{3} \, t)$
}

\begin{exercise}
\pagebreak[2]
Gerbong \unit[5000]{kg} menabrak penahan (pegas) dengan \unitfrac[1]{m}{s}, dan pegas tertekan \unit[0.1]{m}. Andaikan tanpa redaman.
\begin{tasks}
\task Temukan $k$.
\task Seberapa jauh pegas tertekan ketika gerbong \unit[10000]{kg} menabraknya dengan kecepatan yang sama?
\task Jika pegas patah apabila tertekan lebih dari \unit[0.3]{m}, berapa massa maksimum gerbong yang dapat menabraknya dengan \unitfrac[1]{m}{s}?
\task Berapa massa maksimum gerbong yang dapat menabrak pegas tanpa mematahkannya dengan \unitfrac[2]{m}{s}?
\end{tasks}
\end{exercise}
\exsol{%
a) $k=\unitfrac[500000]{N}{m}$
\quad
b) $\frac{1}{5\sqrt{2}} \approx \unit[0.141]{m}$
\quad
c) \unit[45000]{kg}
\quad
d) \unit[11250]{kg}
}

\begin{exercise}
Massa $m$ \unit{kg} berada pada pegas dengan $k=\unitfrac[3]{N}{m}$ dan $c=\unitfrac[2]{Ns}{m}$. Temukan massa $m_0$ yang menghasilkan redaman kritis. Jika $m < m_0$, apakah sistem berosilasi, yakni apakah teredam kurang atau teredam berlebih?
\end{exercise}
\exsol{%
$m_0 = \unit[\nicefrac{1}{3}]{kg}$. Jika $m < m_0$, sistem teredam berlebih dan tidak berosilasi.
}

%%%%%%%%%%%%%%%%%%%%%%%%%%%%%%%%%%%%%%%%%%%%%%%%%%%%%%%%%%%%%%%%%%%%%%%%%%%%%%

\sectionnewpage
\section{Persamaan nonhomogen}
\label{sec:nonhom}

%mbxINTROSUBSECTION

\sectionnotes{2 kuliah\EPref{, \S3.5 dalam \cite{EP}}\BDref{,
\S3.5 dan \S3.6 dalam \cite{BD}}}

\subsection{Menyelesaikan persamaan nonhomogen}

Kita telah menyelesaikan persamaan homogen linear berkoefisien konstan. Bagaimana dengan ODE linear nonhomogen? Misalnya, persamaan getaran mekanis yang dipaksa.
%Now suppose that we drop the requirement of homogeneity.
%This
%usually corresponds to some outside input to the system we are trying to
%model, like the forcing function for the mechanical vibrations of last
%section.
Perhatikan persamaan seperti
\begin{equation} \label{eq3.5:nh}
y'' + 5y'+ 6y = 2x+1 .
\end{equation}
%We still say this equation is a constant-coefficient equation.  We
%only require constants in front of the $y''$, $y'$, and $y$.
%
Kita akan menulis $Ly = 2x+1$ ketika bentuk persis operator tidak penting. Kita menyelesaikan \eqref{eq3.5:nh} sebagai berikut. Pertama, temukan solusi umum $y_c$ untuk \emph{\myindex{persamaan homogen terkait}}
\begin{equation} \label{eq3.5:h}
y'' + 5y'+ 6y = 0 .
\end{equation}
Kita menyebut $y_c$ sebagai \emph{\myindex{solusi komplementer}}. Berikutnya, dengan suatu cara temukan satu \emph{\myindex{solusi partikular}} $y_p$ dari \eqref{eq3.5:nh}. Maka
\begin{equation*}
y = y_c + y_p
\end{equation*}
adalah solusi umum \eqref{eq3.5:nh}. Kita mempunyai $L y_c = 0$ dan $L y_p = 2x+1$. Karena $L$ adalah \emph{\myindex{operator linear}}, kita verifikasi bahwa $y$ adalah solusi: $L y = L ( y_c + y_p) = L y_c + L y_p = 0 + (2x+1)$.

Mari lihat mengapa kita memperoleh solusi \emph{umum}. Misalkan $y_p$ dan $\tilde{y}_p$ adalah dua solusi partikular berbeda dari \eqref{eq3.5:nh}. Tuliskan selisihnya sebagai $w = y_p - \tilde{y}_p$. Substitusikan $w$ ke ruas kiri untuk memperoleh
\begin{equation*}
w'' + 5w'+ 6w =
(y_p'' + 5y_p'+ 6y_p) -
(\tilde{y}_p'' + 5\tilde{y}_p'+ 6\tilde{y}_p) =
(2x+1) - (2x+1) = 0 .
\end{equation*}
Dengan notasi operator, perhitungannya lebih jelas.
%As $L$ is a \emph{\myindex{linear operator}} and so we could just
Karena $L$ operator linear,
\begin{equation*}
Lw = L(y_p - \tilde{y}_p) =
Ly_p - L\tilde{y}_p =
(2x+1)-(2x+1) = 0 .
\end{equation*}
Jadi $w = y_p - \tilde{y}_p$ adalah solusi \eqref{eq3.5:h}, yakni $Lw = 0$. Setiap dua solusi \eqref{eq3.5:nh} berbeda sebesar suatu solusi persamaan homogen \eqref{eq3.5:h}. Solusi $y = y_c + y_p$ mencakup \emph{semua} solusi \eqref{eq3.5:nh}, karena $y_c$ adalah solusi umum persamaan homogen terkait.

\begin{theorem}
Misalkan $Ly=f(x)$ suatu ODE linear (tidak harus berkoefisien konstan). Misalkan $y_c$ solusi komplementer (solusi umum persamaan homogen terkait $Ly = 0$) dan $y_p$ sebarang solusi partikular $Ly=f(x)$. Maka solusi umum $Ly=f(x)$ adalah
\begin{equation*}
y = y_c + y_p.
\end{equation*}
\end{theorem}

Intinya, kita dapat menemukan solusi partikular dengan cara apa pun. Jika kita menemukan solusi partikular lain (dengan metode lain atau sekadar menebak), solusi umumnya tetap sama. Rumusnya mungkin tampak berbeda dan konstanta yang harus dipilih untuk memenuhi syarat awal mungkin berbeda, tetapi solusinya sama.

\subsection{Koefisien tak tentu}
\index{koefisien tak tentu}

Kuncinya ialah menebak secara cerdas satu solusi partikular \eqref{eq3.5:nh}. Perhatikan bahwa $2x+1$ adalah polinomial, dan ruas kiri persamaan akan berupa polinomial jika $y$ dipilih sebagai polinomial berderajat sama. Coba
\begin{equation*}
y_p = Ax + B .
\end{equation*}
Substitusikan $y_p$ ke ruas kiri untuk memperoleh
\begin{equation*}
\begin{split}
y_p'' + 5y_p'+ 6y_p & =
(Ax+B)'' + 5(Ax+B)' + 6(Ax+B)
\\
& =
0 + 5A + 6Ax + 6B = 6Ax+ (5A+6B) .
\end{split}
\end{equation*}
Jadi $6Ax+(5A+6B) = 2x+1$. Maka $A = \nicefrac{1}{3}$ dan $B = \nicefrac{-1}{9}$, sehingga $y_p = \frac{1}{3}\, x - \frac{1}{9} = \frac{3x-1}{9}$. Dengan menyelesaikan masalah komplementer (latihan!), diperoleh
\begin{equation*}
y_c = C_1 e^{-2x} + C_2 e^{-3x}.
\end{equation*}
Jadi solusi umum \eqref{eq3.5:nh} adalah
\begin{equation*}
y = C_1 e^{-2x} + C_2 e^{-3x} + \frac{3x-1}{9} .
\end{equation*}
Andaikan selanjutnya diberikan syarat awal $y(0) = 0$ dan $y'(0) = \nicefrac{1}{3}$. Pertama temukan $y' = - 2C_1 e^{-2x} - 3C_2 e^{-3x} + \nicefrac{1}{3}$. Maka
\begin{equation*}
0 = y(0) = C_1 + C_2 -\frac{1}{9} , \qquad
\frac{1}{3} = y'(0) = - 2C_1 - 3C_2 + \frac{1}{3} .
\end{equation*}
Penyelesaiannya $C_1 = \nicefrac{1}{3}$ dan $C_2 = \nicefrac{-2}{9}$. Solusi khusus yang diinginkan adalah
\begin{equation*}
y = \frac{1}{3} e^{-2x} - \frac{2}{9} e^{-3x} + \frac{3x-1}{9} =
\frac{3 e^{-2x} - 2 e^{-3x} + 3x-1}{9} .
\end{equation*}

\begin{exercise}
Periksa bahwa $y$ benar-benar menyelesaikan persamaan \eqref{eq3.5:nh} dan syarat awal yang diberikan.
\end{exercise}

Catatan: Kesalahan umum adalah menentukan konstanta menggunakan syarat awal pada $y_c$, lalu baru menambahkan solusi partikular $y_p$. Cara itu \emph{tidak} akan berhasil. Anda harus terlebih dahulu menghitung $y = y_c + y_p$, dan \emph{baru kemudian} menentukan konstanta menggunakan syarat awal.

Catatan penting lain: jangan melupakan suku berderajat lebih rendah, sekalipun tidak muncul di ruas kanan. Jika persamaannya $Ly=x^3+1$, Anda harus mencoba $y_p = Ax^3+Bx^2+Cx+D$, walaupun tidak ada $x^2$ ataupun $x$ di ruas kanan. Yang harus dicoba adalah polinomial umum derajat 3.

\medskip

Ruas kanan yang terdiri dari eksponensial, sinus, dan kosinus dapat ditangani dengan cara serupa. Sebagai contoh,
\begin{equation*}
y''+2y'+2y = \cos (2x) .
\end{equation*}
Mari cari suatu $y_p$. Mulailah dengan menebak bahwa solusi memuat suatu kelipatan $\cos(2x)$. Mungkin perlu ditambahkan pula kelipatan $\sin (2x)$ karena turunan kosinus adalah sinus. Coba
\begin{equation*}
y_p = A \cos (2x) + B \sin (2x) .
\end{equation*}
Substitusi $y_p$ ke persamaan memberikan
\begin{multline*}
\underbrace{-4 A \cos (2x) - 4 B \sin (2x)}_{y_p''}
+2 \underbrace{\bigl(-2A \sin (2x) + 2B \cos (2x)\bigr)}_{y_p'}
\\
+
2 \underbrace{\bigl(A \cos (2x) + B \sin (2x)\bigr)}_{y_p}
= \cos (2x) ,
\end{multline*}
atau
\begin{equation*}
(-4A+4B+2A) \cos(2x) +
(-4B-4A+2B) \sin(2x)
= \cos(2x) .
\end{equation*}
Ruas kiri harus sama dengan ruas kanan. Jadi $-4A + 4B + 2A = 1$ dan $-4B - 4A + 2B = 0$. Maka $-2A+4B =1$ dan $2A+B=0$, sehingga $A=\nicefrac{-1}{10}$ dan $B=\nicefrac{1}{5}$. Jadi
\begin{equation*}
y_p = A \cos (2x) + B \sin (2x) = \frac{-\cos (2x) + 2 \sin (2x)}{10} .
\end{equation*}

Demikian pula, jika ruas kanan memuat eksponensial, kita mencoba eksponensial. Jika
%if the equation is (where $L$ is a linear constant-coefficient operator)
\begin{equation*}
Ly = e^{3x},
\end{equation*}
kita menebak $y_p = A e^{3x}$ lalu mencoba menentukan $A$.

\medskip

Ketika ruas kanan merupakan kelipatan sinus, kosinus, eksponensial, dan polinomial, kita dapat menggunakan aturan hasil kali untuk diferensiasi guna menyusun tebakan. Kita perlu menebak bentuk $y_p$ sedemikian sehingga $Ly_p$ berbentuk sama dan memuat semua suku yang diperlukan ruas kanan. Sebagai contoh,
\begin{equation*}
Ly = (1+3x^2)\,e^{-x}\cos (\pi x) .
\end{equation*}
Untuk persamaan ini, tebak
\begin{equation*}
y_p = (A + Bx + Cx^2)\,e^{-x} \cos (\pi x) +
(D + Ex + Fx^2)\,e^{-x} \sin (\pi x) .
\end{equation*}
Substitusikan, lalu mudah-mudahan diperoleh persamaan yang dapat diselesaikan untuk $A$, $B$, $C$, $D$, $E$, dan $F$. Seperti terlihat, perhitungan ini dapat dengan cepat menjadi sangat panjang dan \myindex{menjemukan}. % a bit of fun
C'est \myindex{la vie}! %bit more fun

\medskip

Ada satu kendala. Tebakan kita mungkin justru menyelesaikan persamaan homogen terkait. Perhatikan
\begin{equation*}
y'' - 9y = e^{3x} .
\end{equation*}
Kita ingin menebak $y_p = Ae^{3x}$, tetapi substitusi ke ruas kiri memberikan
\begin{equation*}
y_p''-9y_p = 9Ae^{3x} - 9Ae^{3x} = 0 \neq e^{3x} .
\end{equation*}
Tidak ada pilihan $A$ yang membuat ruas kiri menjadi $e^{3x}$. Kuncinya adalah mengalikan tebakan dengan $x$ untuk menghilangkan duplikasi dengan solusi komplementer. Pertama, temukan $y_c$ (solusi $Ly = 0$),
\begin{equation*}
y_c = C_1 e^{-3x} + C_2 e^{3x} .
\end{equation*}
Perhatikan bahwa suku $e^{3x}$ menduplikasi tebakan yang diinginkan. Ubah tebakan menjadi $y_p = Axe^{3x}$ agar tidak ada duplikasi. Coba: $y_p' = Ae^{3x} + 3Axe^{3x}$ dan $y_p'' = 6Ae^{3x} + 9Axe^{3x}$, sehingga
\begin{equation*}
y_p'' -9y_p = 6Ae^{3x} + 9Axe^{3x} - 9Axe^{3x} =
6Ae^{3x} .
\end{equation*}
Jadi $6Ae^{3x}$ harus sama dengan $e^{3x}$. Maka $6A = 1$ dan $A=\nicefrac{1}{6}$. Kini solusi umum dapat ditulis
\begin{equation*}
y = y_c + y_p =
C_1 e^{-3x} + C_2 e^{3x} + \frac{1}{6}\,xe^{3x} .
\end{equation*}

\medskip

Mungkin saja perkalian dengan $x$ tidak menghilangkan semua duplikasi. Perhatikan
\begin{equation*}
y''-6y'+9y = e^{3x} .
\end{equation*}
Solusi komplementernya $y_c = C_1 e^{3x} + C_2 x e^{3x}$. Menebak $y_p=A xe^{3x}$ tidak membantu. Kita harus menebak $y_p = Ax^2e^{3x}$. Pada dasarnya, kalikan tebakan dengan $x$ sampai semua duplikasi hilang. \emph{Tetapi jangan lebih dari itu!} Mengalikan terlalu banyak tidak akan berhasil.

Pastikan pula mencari duplikasi dalam seluruh ekspresi yang dicoba. Untuk persamaan $y''-y=xe^x$, dengan $y_c = C_1e^x + C_2e^{-x}$, tidak cukup bahwa $xe^x$ tidak muncul dalam $y_c$. Mencoba $y_p=Ae^x+Bxe^x$ tidak akan berhasil karena $e^x$ muncul dalam $y_c$; Anda harus mencoba $y_p=Axe^x+Bx^2e^x$.

\medskip

Terakhir, bagaimana jika ruas kanan mempunyai beberapa suku? Misalnya,
\begin{equation*}
Ly = e^{2x} + \cos x .
\end{equation*}
Dalam kasus ini, temukan $u$ yang menyelesaikan $Lu = e^{2x}$ dan $v$ yang menyelesaikan $Lv = \cos x$ (kerjakan setiap suku secara terpisah). Jika $y_p = u+ v$, diperoleh solusi partikular yang diinginkan. Alasannya $L$ linear: $Ly_p = L(u+v) = Lu + Lv = e^{2x} + \cos x$.

\subsection{Variasi parameter}

Metode koefisien tak tentu berlaku untuk banyak masalah dasar, tetapi tidak selalu. Metode itu hanya berlaku ketika ruas kanan $Ly = f(x)$ mempunyai berhingga banyak turunan yang bebas linear, sehingga kita dapat menulis tebakan yang mencakup semuanya. Beberapa persamaan lebih sulit. Perhatikan
\begin{equation*}
y''+y = \tan x .
\end{equation*}
Setiap turunan baru $\tan x$ tampak sama sekali berbeda dan tidak dapat ditulis sebagai kombinasi linear turunan sebelumnya. Jika kita mulai mendiferensialkan $\tan x$, diperoleh:
\begin{multline*}
\sec^2 x, \quad
2\sec^2 x \, \tan x, \quad
4 \sec^2 x \, \tan^2 x + 2 \sec^4 x, \\
8 \sec^2 x \, \tan^3 x + 16 \sec^4 x \, \tan x, \quad
16\sec^2 x \, \tan^4 x + 88 \sec^4 x \tan^2 x + 16 \sec^6 x, \quad
\ldots
\end{multline*}

Persamaan ini memerlukan metode lain. Kita sajikan metode \emph{\myindex{variasi parameter}}, yang menangani setiap persamaan berbentuk $Ly = f(x)$ asalkan integral tertentu dapat diselesaikan. Demi kesederhanaan, kita batasi pada persamaan orde kedua berkoefisien konstan, tetapi metode ini juga berlaku untuk persamaan orde tinggi (perhitungannya menjadi lebih \myindex{menjemukan}). % a bit of fun
Metode ini juga berlaku untuk persamaan berkoefisien tak konstan, asalkan persamaan homogen terkait dapat diselesaikan.

Rincian berikut berlaku untuk sebarang persamaan berbentuk $y''+p(x)y' +q(x)y = f(x)$, tetapi mungkin lebih baik menjelaskan metode melalui contoh khusus. Perhatikan
\begin{equation*}
Ly = y''+y = \tan x .
\end{equation*}
Pertama temukan solusi komplementer (solusi $Ly_c = 0$). Diperoleh $y_c = C_1 y_1 + C_2 y_2$, dengan $y_1 = \cos x$ dan $y_2 = \sin x$. Untuk menemukan solusi partikular persamaan nonhomogen, kunci metode ini adalah mencoba
\begin{equation*}
y_p = y = u_1 y_1 + u_2 y_2 ,
\end{equation*}
di mana $u_1$ dan $u_2$ adalah \emph{fungsi}, bukan konstanta. Kita mencoba memenuhi $Ly = \tan x$. Ini memberikan satu syarat pada fungsi $u_1$ dan $u_2$. Hitung (perhatikan aturan hasil kali!)
\begin{equation*}
y' = (u_1' y_1 + u_2' y_2) + (u_1 y_1' + u_2 y_2').
\end{equation*}
Kita masih dapat menetapkan satu syarat tambahan untuk menyederhanakan perhitungan (ada dua fungsi tak diketahui, sehingga dua syarat semestinya diperbolehkan). Syaratkan $u_1' y_1 + u_2' y_2 = 0$ (suku pertama di atas). Ini mempermudah perhitungan turunan kedua:
\begin{align*}
& y' = u_1 y_1' + u_2 y_2' , \\
& y'' = (u_1' y_1' + u_2' y_2') + (u_1 y_1'' + u_2 y_2'') .
\end{align*}
Karena $y_1$ dan $y_2$ adalah solusi $y''+y = 0$, diperoleh $y_1'' = - y_1$ dan $y_2'' = - y_2$. (Jika persamaan lebih umum $y''+p(x)y' +q(x)y = 0$, kita akan mempunyai $y_i'' = -p(x)y_i' -q(x)y_i$.) Jadi
\begin{equation*}
y'' = (u_1' y_1' + u_2' y_2') - (u_1 y_1 + u_2 y_2) .
\end{equation*}
Karena $u_1 y_1 + u_2 y_2 = y$,
\begin{equation*}
y'' = (u_1' y_1' + u_2' y_2') - y ,
\end{equation*}
dan karenanya
\begin{equation*}
y'' + y = Ly = u_1' y_1' + u_2' y_2' .
\avoidbreak
\end{equation*}
Supaya $y$ memenuhi $Ly = f(x)$, harus berlaku $f(x) = u_1' y_1' + u_2' y_2'$.

Yang perlu diselesaikan adalah dua persamaan (syarat) yang kita tetapkan pada $u_1$ dan $u_2$:
\begin{equation*}
\mybxbg{~~
\begin{aligned}
& u_1' y_1 + u_2' y_2 = 0 ,\\
& u_1' y_1' + u_2' y_2' = f(x) .
\end{aligned}
~~}
\end{equation*}
Persamaan yang sama diperoleh untuk sebarang $Ly = f(x)$, dengan $Ly = y''+p(x)y'+q(x)y$. Kita menentukan $u_1'$ dan $u_2'$ dalam $f(x)$, $y_1$, dan $y_2$. Ada rumus umum yang dapat langsung digunakan, tetapi daripada menghafalnya, lebih mudah menyelesaikannya seperti di bawah. Dalam kasus kita, $y_1 = \cos x$, $y_2 = \sin x$, dan $f(x) = \tan x$, sehingga kedua persamaan adalah
\begin{equation*}
\begin{aligned}
u_1' \cos x + u_2' \sin x &= 0 ,\\
-u_1' \sin x + u_2' \cos x &= \tan x .
\end{aligned}
\end{equation*}
Kalikan persamaan pertama dengan $\sin x$ dan yang kedua dengan $\cos x$:
\begin{equation*}
\begin{aligned}
u_1' \cos x \sin x + u_2' \sin^2 x & = 0 ,\\
-u_1' \sin x \cos x + u_2' \cos^2 x & = \tan x \cos x = \sin x .
\end{aligned}
\end{equation*}
Jumlahkan kedua persamaan untuk mengeliminasi $u_1'$, tentukan $u_2'$, lalu $u_1'$:
\begin{equation*}
\begin{aligned}
& u_2' \bigl(\sin^2 x + \cos^2 x\bigr) = \sin x , \\
& u_2' = \sin x , \\
& u_1' = \frac{- \sin^2 x}{\cos x} = \cos x - \sec x .
\end{aligned}
\end{equation*}
Integralkan $u_1'$ dan $u_2'$ untuk memperoleh $u_1$ dan $u_2$.
\begin{equation*}
\begin{aligned}
& u_1 = \int u_1'\,dx
= \int ( \cos x-\sec x ) \,dx
= \sin x -
\ln \left\lvert \sec x + \tan x \right\rvert
, \\
& u_2 = \int u_2'\,dx
= \int \sin x\,dx = -\cos x .
\end{aligned}
\end{equation*}
Kita mencari solusi partikular, jadi konstanta integrasi diabaikan. Solusi partikularnya
\begin{multline*}
y_p = u_1 y_1 + u_2 y_2 =
\cos x \sin x
-
\cos x
\,
\ln \lvert
\sec x + \tan x
\rvert
-\cos x \sin x
\\ =
-
\cos x \, \ln \lvert
\sec x + \tan x
\rvert .
\end{multline*}
Jadi solusi umum $y'' + y = \tan x$ adalah
\begin{equation*}
y = C_1 \cos x + C_2 \sin x
-
\cos x \, \ln \lvert
\sec x + \tan x
\rvert .
\end{equation*}

\medskip

Gagasan umum untuk sebarang $y'' + p y' + q y = f(x)$ adalah pertama-tama menemukan solusi $y_1$, $y_2$ dari $y''+py'+qy = 0$. Lalu selesaikan dua persamaan berbingkai untuk $u_1'$ dan $u_2'$, yakni $u_1' y_1 + u_2' y_2 = 0$ dan $u_1' y_1' + u_2' y_2' = f(x)$. Integralkan $u_1'$ dan $u_2'$ untuk menemukan $u_1$ dan $u_2$, lalu substitusikan ke $y = u_1 y_1 + u_2 y_2$ untuk menemukan solusi partikular. Perhatikan bahwa jika $y''$ mempunyai koefisien selain $1$, yakni persamaannya $ay'' + by' + cy = f(x)$, Anda harus lebih dahulu membagi persamaan (dan karenanya $f$) dengan $a$.

\subsection{Latihan}

\begin{exercise}
Temukan solusi partikular $y''-y' -6y = e^{2x}$.
\end{exercise}

\begin{exercise}
Temukan solusi partikular $y''-4y' +4y = e^{2x}$.
\end{exercise}

\begin{exercise}
Selesaikan masalah nilai awal $y''+9y = \cos (3x) + \sin (3x)$ untuk $y(0) = 2$, $y'(0) = 1$.
\end{exercise}

\begin{exercise}
Susun bentuk solusi partikular, tetapi jangan tentukan koefisiennya, untuk $y^{(4)}-2y'''+y'' = e^x$.
\end{exercise}

\begin{exercise}
Susun bentuk solusi partikular, tetapi jangan tentukan koefisiennya, untuk $y^{(4)}-2y'''+y'' = e^x + x + \sin x$.
\end{exercise}

\begin{exercise}
\pagebreak[2]
\leavevmode
\begin{tasks}
\task Dengan variasi parameter, temukan solusi partikular $y''-2y'+y = e^x$.
\task Temukan solusi partikular menggunakan koefisien tak tentu.
\task Apakah kedua solusi yang ditemukan sama? Lihat juga \exerciseref{exercise:diffvarparunder}.
\end{tasks}
\end{exercise}

\begin{exercise}
Temukan solusi partikular $y''-2y' +y = \sin (x^2)$. Jawaban boleh dibiarkan sebagai integral tentu.
\end{exercise}

\begin{exercise}
Untuk konstanta sebarang $c$, temukan solusi partikular $y''-y=e^{cx}$. Petunjuk: Pastikan menangani setiap kemungkinan $c$ real.
\end{exercise}

\begin{exercise} \label{exercise:diffvarparunder}
\pagebreak[2]
\leavevmode
\begin{tasks}
\task Dengan variasi parameter, temukan solusi partikular $y''-y = e^x$.
\task Temukan solusi partikular menggunakan koefisien tak tentu.
\task Apakah kedua solusi yang ditemukan sama? Apa yang terjadi?
\end{tasks}
\end{exercise}

\begin{exercise}
Temukan polinomial $P(x)$ sedemikian sehingga $y = 2 x^2 + 3 x + 4$ menyelesaikan $y''+5 y'+ y = P(x)$.
\end{exercise}

\setcounter{exercise}{100}

\begin{exercise}
Temukan solusi partikular $y''-y'+y=2\sin(3x)$.
\end{exercise}
\exsol{%
$y=\frac{-16\sin(3x)+6\cos(3x)}{73}$
}

\begin{samepage}
\begin{exercise}
\leavevmode
\begin{tasks}
\task Temukan solusi partikular $y''+2y=e^x + x^3$.
\task Temukan solusi umum.
\end{tasks}
\end{exercise}
\end{samepage}
\exsol{%
a) $y=\frac{2e^x+3x^3-9x}{6}$ \quad
b) $y=C_1 \cos(\sqrt{2} x) + C_2 \sin(\sqrt{2} x) + \frac{2e^x+3x^3-9x}{6}$
}

\begin{exercise}
Selesaikan $y''+2y'+y = x^2$, $y(0)=1$, $y'(0)=2$.
\end{exercise}
\exsol{%
$y(x) = x^2-4 x+6+e^{-x}(x-5)$
}

\begin{exercise}
Gunakan variasi parameter untuk menemukan solusi partikular $y''-y = \frac{1}{e^x+e^{-x}}$.
\end{exercise}
\exsol{%
$y = \frac{2xe^x-(e^x+e^{-x})\ln(e^{2x}+1)}{4}$
}

\begin{exercise}
Untuk konstanta sebarang $c$, temukan solusi umum $y''-2y=\sin(x+c)$.
\end{exercise}
\exsol{%
$y=\frac{-\sin(x+c)}{3}+C_1 e^{\sqrt{2}\,x}+C_2 e^{-\sqrt{2}\,x}$
}

\begin{exercise}
Koefisien tak tentu kadang dapat digunakan untuk menebak solusi partikular persamaan selain yang berkoefisien konstan. Temukan polinomial $y(x)$ yang menyelesaikan $y'+ x y = x^3+2x^2+5x+2$.

Catatan: Tidak setiap ruas kanan memungkinkan solusi polinomial; misalnya $y'+xy=1$ tidak, tetapi teknik berdasarkan koefisien tak tentu dapat digunakan, lihat \chapterref{ps:chapter}.
\end{exercise}
\exsol{%
$y=x^2+2x+3$
}


%%%%%%%%%%%%%%%%%%%%%%%%%%%%%%%%%%%%%%%%%%%%%%%%%%%%%%%%%%%%%%%%%%%%%%%%%%%%%%

\sectionnewpage
\section{Osilasi dipaksa dan resonansi} \label{forcedo:section}

%mbxINTROSUBSECTION

\sectionnotes{2 kuliah\EPref{, \S3.6 dalam \cite{EP}}\BDref{,
\S3.8 dalam \cite{BD}}}

\begin{mywrapfigsimp}{2.0in}{2.3in}
\noindent
\inputpdft{massfigforce}{%
%note that this diagram appears more than once in the book
Diagram yang menunjukkan massa berlabel m, terhubung di sebelah kiri ke pegas berlabel k yang terhubung ke dinding di sebelah kiri. Panah gaya ke kanan pada massa berlabel F dari t. Panah di bawah massa yang menunjukkan perpindahan mengarah ke kanan dan berlabel x. `Lantai' diagram berlabel `redaman c'.}
\end{mywrapfigsimp}
Mari kembali ke contoh massa pada pegas. Kita meninjau kasus osilasi dipaksa yang belum ditangani. Jadi kita perhatikan persamaan
\begin{equation*}
mx'' + cx' + kx = F(t) ,
\end{equation*}
untuk suatu $F(t)$ taknol. Susunannya kembali sebagai berikut: $x$ adalah posisi, $m$ massa, $c$ gesekan, $k$ konstanta pegas, dan $F(t)$ gaya eksternal yang bekerja pada massa.

Kita tertarik pada pemaksaan periodik, seperti komponen berputar yang tidak terpusat, suara keras, atau sumber gaya periodik lain. Dengan menggunakan deret Fourier, lihat \chapterref{FS:chapter}, semua fungsi periodik dapat dipahami dengan meninjau $F(t) = F_0 \cos (\omega t)$ (atau sinus sebagai ganti kosinus; perhitungannya pada dasarnya sama), sehingga kita berfokus pada kasus sederhana ini.

\subsection{Gerak dipaksa tak teredam dan resonansi}

Pertama, perhatikan gerak tak teredam ($c=0$). Persamaannya
\begin{equation*}
mx'' + kx = F_0 \cos (\omega t) .
\end{equation*}
Persamaan ini mempunyai solusi komplementer (solusi persamaan homogen terkait)
\begin{equation*}
x_c = C_1 \cos (\omega_0 t) + C_2 \sin (\omega_0 t) ,
\end{equation*}
di mana $\omega_0 = \sqrt{\nicefrac{k}{m}}$ adalah \emph{\myindex{frekuensi alami}} (sudut). Inilah frekuensi ketika sistem \myquote{ingin berosilasi} tanpa gangguan eksternal.

Andaikan $\omega_0 \neq \omega$. Kita menyelesaikannya dengan metode koefisien tak tentu. Coba solusi $x_p = A \cos (\omega t)$ dan tentukan $A$. Kita tidak memerlukan sinus dalam solusi coba karena setelah substitusi hanya ada kosinus. Menyertakan sinus tidak masalah; koefisiennya akan ditemukan nol (saya tidak menemukan rima kedua). Substitusi ke persamaan dan penyelesaian untuk $A$ memberikan
\begin{equation*}
x_p = \frac{F_0}{m(\omega_0^2 - \omega^2)} \cos (\omega t) .
\end{equation*}
Aljabar yang diperlukan diserahkan sebagai latihan.

Solusi umumnya
\begin{equation*}
\mybxbg{
~~
x = C_1 \cos (\omega_0 t) + C_2 \sin (\omega_0 t) +
\frac{F_0}{m(\omega_0^2 - \omega^2)} \cos (\omega t) .
~~
}
\end{equation*}
Dalam bentuk lain,
\begin{equation*}
x = C \cos (\omega_0 t - \gamma) +
\frac{F_0}{m(\omega_0^2 - \omega^2)} \cos (\omega t) .
\end{equation*}
Solusi merupakan superposisi dua gelombang kosinus (tergeser) dengan frekuensi berbeda.
\pagebreak[2]

\begin{example}
Ambil
\begin{equation*}
0.5 x'' + 8 x = 10 \cos (\pi t), \qquad x(0)=0, \qquad x'(0)=0 .
\end{equation*}

Mari kita hitung. Pertama, baca parameternya: $\omega = \pi$, $\omega_0 = \sqrt{\nicefrac{8}{0.5}} = 4$, $F_0 = 10$, $m=0.5$. Solusi umumnya
\begin{equation*}
x = C_1 \cos (4 t) + C_2 \sin (4 t) +
\frac{20}{16 - \pi^2} \cos (\pi t) .
\end{equation*}

%15 is the number of lines, must be adjusted
%mbxSTARTIGNORE
\begin{mywrapfig}[15]{3.25in}
\capstart
\diffyincludegraphics{width=3in}{width=4.5in}{3-6-beating}{%
Grafik kurva yang berosilasi cepat dengan amplitudo bermula dari nol, lalu membesar, kembali ke nol, membesar lagi, dan seterusnya.}
\caption{Grafik
$\frac{20}{16 - \pi^2} \bigl( \cos (\pi t)- \cos (4 t) \bigr)$.\label{3.6:beatingfig}}
\end{mywrapfig}
%mbxENDIGNORE
%
% make sure the MBX below is synced!
%

Tentukan $C_1$ dan $C_2$ menggunakan syarat awal: $C_1 = \frac{-20}{16 - \pi^2}$ dan $C_2 = 0$. Jadi
\begin{equation*}
x =
\frac{20}{16 - \pi^2} \bigl( \cos (\pi t)- \cos (4 t) \bigr) .
\end{equation*}

%mbxlatex \begin{myfig}
%mbxlatex \capstart
%mbxlatex \diffyincludegraphics{width=3in}{width=4.5in}{3-6-beating}{%
%mbxlatex A plot of a quickly oscillating curve whose amplitude starts at zero, then
%mbxlatex grows, then goes to zero, then grows again, etc.}
%mbxlatex \caption{Graph of
%mbxlatex $\frac{20}{16 - \pi^2} \bigl( \cos (\pi t)- \cos (4 t) \bigr)$.\label{3.6:beatingfig}}
%mbxlatex \end{myfig}

Perhatikan perilaku \myquote{layangan}\index{layangan} dalam \figurevref{3.6:beatingfig}. Untuk menganalisisnya, gunakan identitas trigonometri
\begin{equation*}
2\sin \left( \frac{A-B}{2} \right) \sin \left( \frac{A+B}{2} \right) =
\cos B -\cos A
\end{equation*}
untuk memperoleh
\begin{equation*}
x =
\frac{20}{16 - \pi^2} \left( 2 \sin \left(\frac{4-\pi}{2} t \right)
\sin \left( \frac{4+\pi}{2} t \right) \right) .
\end{equation*}
Fungsi $x$ adalah gelombang berfrekuensi tinggi yang dimodulasi oleh gelombang berfrekuensi rendah.
\end{example}

Sekarang andaikan $\omega_0 = \omega$. Perhatikan bahwa $\cos (\omega t)$ menyelesaikan persamaan homogen terkait. Jadi kita tidak dapat mencoba solusi $A \cos (\omega t)$ dengan metode koefisien tak tentu. Karena itu, coba $x_p = A t \cos (\omega t) + B t \sin (\omega t)$. Kali ini suku sinus diperlukan karena turunan kedua $t \cos (\omega t)$ memuat sinus. Tuliskan persamaan
\begin{equation*}
x'' + \omega^2 x = \frac{F_0}{m} \cos ( \omega t) .
\end{equation*}
Substitusi $x_p$ ke ruas kiri memberikan
\begin{equation*}
2 B \omega \cos (\omega t) - 2 A \omega \sin (\omega t) =
\frac{F_0}{m} \cos (\omega t) .
\end{equation*}
Jadi $A = 0$ dan $B = \frac{F_0}{2m\omega}$. Solusi partikularnya $\frac{F_0}{2m\omega} \, t \sin (\omega t)$ dan solusi umumnya
\begin{equation*}
x = C_1 \cos (\omega t) + C_2 \sin (\omega t)
+ \frac{F_0}{2m\omega} \, t \sin (\omega t) .
\end{equation*}

Suku penting adalah suku terakhir (solusi partikular). Suku ini tumbuh tanpa batas saat $t \to \infty$. Bahkan, ia berosilasi antara $\frac{F_0 t}{2m\omega}$ dan $\frac{- F_0 t}{2m\omega}$. Dua suku pertama hanya berosilasi antara $\pm\sqrt{C_1^2 + C_2^2}$, yang secara proporsional semakin kecil dibandingkan osilasi suku terakhir ketika $t$ membesar. Dalam \figurevref{3.6:resonancefig}, terlihat grafik dengan $C_1=C_2=0$, $F_0 = 2$, $m=1$, $\omega = \pi$.

\begin{mywrapfig}{3.25in}
\capstart
\diffyincludegraphics{width=3in}{width=4.5in}{3-6-resonance}{%
Kurva yang berosilasi cepat dengan amplitudo bermula dari nol, lalu bertumbuh secara linear ketika bergerak ke kanan dalam grafik.}
\caption{Grafik
$\frac{1}{\pi} t \sin (\pi t)$.\label{3.6:resonancefig}}
\end{mywrapfig}

Dengan memaksa sistem tepat pada frekuensi yang sesuai, kita menghasilkan osilasi yang sangat liar. Perilaku ini disebut \emph{\myindex{resonansi}} atau mungkin \emph{\myindex{resonansi murni}}. Kadang resonansi diinginkan. Ingatlah ketika masih kecil Anda dapat mulai berayun hanya dengan bergerak maju-mundur pada kursi ayunan dengan \myquote{frekuensi yang tepat}? Anda sedang mencoba mencapai resonansi. Gaya setiap gerakan kecil, tetapi setelah beberapa saat menghasilkan ayunan besar.

Di sisi lain, resonansi dapat merusak. Dalam gempa bumi, sebagian bangunan runtuh sementara yang lain mungkin relatif tidak rusak. Ini disebabkan bangunan berbeda mempunyai frekuensi resonansi berbeda. Jadi menentukan frekuensi resonansi dapat sangat penting.

Contoh umum (tetapi keliru) mengenai daya rusak resonansi adalah runtuhnya jembatan Tacoma Narrows. Ternyata ada fenomena lain yang berperan%
\footnote{K.\ Billah dan R.\ Scanlan, \emph{Resonansi, Kegagalan Jembatan Tacoma Narrows, dan Buku Teks Fisika Sarjana}, American Journal of Physics, 59(2), 1991, 118--124,
\url{http://www.ketchum.org/billah/Billah-Scanlan.pdf}}.

\subsection{Gerak dipaksa teredam dan resonansi praktis}

Dalam kenyataan, keadaan tidak sesederhana di atas. Tentu terdapat redaman. Persamaannya menjadi
\begin{equation} \label{3.6:deq}
mx'' + cx' + kx = F_0 \cos (\omega t) ,
\end{equation}
untuk suatu $c > 0$. Masalah homogen telah diselesaikan. Tetapkan
\begin{equation*}
p = \frac{c}{2m},  \qquad \omega_0 = \sqrt{\frac{k}{m}} .
\end{equation*}
Ganti persamaan \eqref{3.6:deq} dengan
\begin{equation*}
x'' + 2px' + \omega_0^2x = \frac{F_0}{m} \cos (\omega t) .
\end{equation*}
Akar persamaan karakteristik masalah homogen terkait adalah $r_1,r_2 = -p \pm \sqrt{p^2 - \omega_0^2}$. Bentuk solusi umum persamaan homogen terkait bergantung pada tanda $p^2 - \omega_0^2$, atau secara ekuivalen tanda $c^2 - 4km$, seperti sebelumnya:
\begin{equation*}
x_c =
\begin{cases}
C_1 e^{r_1 t} + C_2 e^{r_2 t} & \text{jika } \; c^2 > 4km , \\
C_1 e^{-p t} + C_2 t e^{-p t} & \text{jika } \; c^2 = 4km , \\
e^{-p t} \bigl( C_1 \cos (\omega_1 t) + C_2 \sin (\omega_1 t) \bigr) &
  \text{jika } \; c^2 < 4km ,
\end{cases}
\end{equation*}
di mana $\omega_1 = \sqrt{\omega_0^2 - p^2}$. Dalam semua kasus, $x_c(t) \to 0$ ketika $t \to \infty$.

\pagebreak[2]
Mari temukan solusi partikular. Tidak mungkin ada konflik ketika menentukan koefisien tak tentu dengan mencoba $x_p = A \cos (\omega t) + B \sin (\omega t)$.
%Hence, we will never get the kind of catastrophic scenario we have seen
%before.
%A slightly different notion of \myquote{resonance} will still occur.
Substitusikan dan tentukan $A$ serta $B$. Diperoleh (rincian \myindex{menjemukan} % a bit of fun
diserahkan kepada pembaca)
\begin{equation*}
\bigl((\omega_0^2  - \omega^2)B - 2\omega p A\bigr) \sin (\omega t)
+
\bigl((\omega_0^2  - \omega^2)A + 2\omega p B\bigr) \cos (\omega t)
=
\frac{F_0}{m} \cos (\omega t) .
\end{equation*}

Kita menentukan $A$ dan $B$ seperti biasa:
\begin{equation*}
A=\frac{(\omega_0^2-\omega^2) F_0}
{m{(2\omega p)}^2+m{(\omega_0^2-\omega^2)}^2} ,
\qquad B=\frac{2 \omega p F_0}
{m{(2\omega p)}^2+m{(\omega_0^2-\omega^2)}^2} .
\end{equation*}
Jadi solusi partikularnya
\begin{equation*}
x_p =
\frac{(\omega_0^2-\omega^2) F_0}
{m{(2\omega p)}^2+m{(\omega_0^2-\omega^2)}^2} \cos (\omega t) +
\frac{2 \omega p F_0}
{m{(2\omega p)}^2+m{(\omega_0^2-\omega^2)}^2} \sin (\omega t) .
\end{equation*}
Amplitudo $C = \sqrt{A^2+B^2}$ dari solusi ini adalah
\begin{equation*}
C = \frac{F_0}{m \sqrt{{(2\omega p)}^2+{(\omega_0^2-\omega^2)}^2}} .
\end{equation*}
Tuliskan solusi dalam bentuk alternatif $x_p = C \cos(\omega t - \gamma)$ dengan amplitudo $C$ dan pergeseran fase $\gamma$, dengan (jika $\omega \neq \omega_0$)
\begin{equation*}
\tan \gamma = \frac{B}{A} = \frac{2\omega p}{\omega_0^2-\omega^2} .
\end{equation*}
Jadi,
\begin{equation*}
\mybxbg{~~
x_p =
\frac{F_0}{m \sqrt{{(2\omega p)}^2+{(\omega_0^2-\omega^2)}^2}}
\cos ( \omega t - \gamma ) .
~~}
\end{equation*}
Jika $\omega = \omega_0$, maka $A=0$, $B = C = \frac{F_0}{2m\omega p}$, dan $\gamma = \nicefrac{\pi}{2}$.

%What is important for us is how this
%solution depends on the parameters, $F_0$, $m$, $\omega$, $\omega_0$,
%and $p$.

%The exact formula is not as important as the idea.  Do not memorize
%the formula above, you should instead remember the ideas involved.
%For a different forcing function $F$, you will get a different formula
%for $x_p$.
%So there is no point in memorizing this specific
%formula.  You can always recompute it later or look it up if you really need
%it.

\medskip

Untuk alasan yang segera dijelaskan, kita menyebut $x_c$ sebagai \emph{\myindex{solusi transien}} dan menyatakannya dengan $x_{tr}$. Kita menyebut $x_p$ sebagai \emph{\myindex{solusi periodik tunak}} dan menyatakannya dengan $x_{sp}$. Solusi umumnya
\begin{equation*}
x = x_c + x_p = x_{tr} + x_{sp} .
\end{equation*}

Solusi transien $x_{tr} = x_c$ menuju nol saat $t \to \infty$, karena semua suku memuat eksponensial berpangkat negatif. Jadi untuk $t$ besar, pengaruh $x_{tr}$ dapat diabaikan dan pada dasarnya hanya $x_{sp}$ yang tampak. Karena itu dinamakan \emph{transien}. Perhatikan bahwa $x_{sp}$ tidak memuat konstanta sebarang, dan syarat awal hanya memengaruhi $x_{tr}$. Jadi pengaruh syarat awal dapat diabaikan setelah beberapa waktu.
%Because of this behavior,
Kita dapat berfokus pada solusi periodik tunak dan mengabaikan solusi transien. Lihat \figurevref{3.6:transbehfig} untuk grafik dengan beberapa syarat awal berbeda.

%15 is the number of lines, must be adjusted
\begin{mywrapfig}[15]{3.25in}
\capstart
\diffyincludegraphics{width=3in}{width=4.5in}{3-6-transbeh}{%
Grafik beberapa solusi yang bermula pada titik berbeda di sisi kiri. Saat t bergerak dari 0 ke 20, semuanya mendekati dan hampir tidak dapat dibedakan dari kurva sinusoidal yang berosilasi.}
\caption{Solusi dengan syarat awal berbeda untuk parameter
$k=1$, $m=1$, $F_0 = 1$, $c=0.7$, dan $\omega=1.1$.\label{3.6:transbehfig}}
\end{mywrapfig}

Seberapa cepat $x_{tr}$ menuju nol bergantung pada $p$ (dan karenanya $c$). Semakin besar $p$ (semakin besar $c$), semakin \myquote{cepat} $x_{tr}$ dapat diabaikan. Sebaliknya, semakin kecil redaman, semakin panjang \myquote{daerah transien.} Ini konsisten dengan pengamatan bahwa ketika $c=0$, syarat awal memengaruhi perilaku untuk sepanjang waktu (yakni \myquote{daerah transien} tak hingga).

\medskip

Mari jelaskan arti resonansi ketika terdapat redaman. Tidak ada konflik ketika menyelesaikan dengan koefisien tak tentu, sehingga tidak ada suku dalam $x_{sp}$ yang menuju tak hingga. Sebagai gantinya, kita meninjau amplitudo $C$ dari solusi periodik tunak $x_{sp}$, yakni jarak terjauh massa ke kedua arah dari posisi diam. Gambarkan $C$ sebagai fungsi $\omega$ (dengan semua parameter lain tetap) dan cari maksimumnya. Nilai $\omega$ yang memberikan maksimum ini adalah \emph{\myindex{frekuensi resonansi praktis}}. Amplitudo maksimum $C(\omega)$ adalah \emph{\myindex{amplitudo resonansi praktis}}. Jadi ketika terdapat redaman, kita berbicara tentang \emph{\myindex{resonansi praktis}}, bukan resonansi murni. \figurevref{3.6:pracresfig} menunjukkan contoh grafik untuk tiga nilai $c$ berbeda. Perhatikan bahwa amplitudo resonansi praktis membesar ketika redaman mengecil, dan resonansi praktis dapat hilang sama sekali ketika redaman besar.

\begin{myfig}
\capstart
\diffyincludegraphics{width=3in}{width=4.5in}{3-6-pracres}{%
Grafik tiga versi fungsi C dari omega untuk omega antara 0 dan 3. Garis atas bermula di 1, naik ke puncak sedikit lebih dari 2.5 tepat sebelum omega sama dengan 1, lalu meluruh menuju nol. Garis tengah kira-kira sama tetapi mencapai puncak yang lebih rendah dan sedikit lebih awal. Garis bawah hanya menurun tanpa mencapai puncak. Ketiganya mempunyai nilai kecil yang kira-kira sama ketika omega sama dengan 3.}
\caption{Grafik $C(\omega)$ yang menunjukkan resonansi praktis dengan parameter
$k=1$, $m=1$, $F_0 = 1$. Garis atas untuk $c=0.4$, garis tengah untuk
$c=0.8$, dan garis bawah untuk
$c=1.6$.\label{3.6:pracresfig}}
\end{myfig}

Untuk menemukan maksimum $C(\omega)$, cari turunan $C'(\omega)$:
\begin{equation*}
C'(\omega) =
\frac{- 2\omega( 2p^2+\omega^2-\omega_0^2)F_0}
{m {\bigl({(2\omega p)}^2+{(\omega_0^2-\omega^2)}^2\bigr)}^{3/2}} .
\end{equation*}
Nilainya nol ketika $\omega = 0$ atau $2p^2+\omega^2-\omega_0^2 = 0$. Dengan kata lain, $C'(\omega) = 0$ ketika
\begin{equation*}
\mybxbg{
~~
\omega = \sqrt{\omega_0^2 - 2p^2} \quad \text{atau} \quad \omega = 0 .
~~
}
\end{equation*}
Jika $\omega_0^2 - 2p^2$ positif, maka $\sqrt{\omega_0^2 - 2p^2}$ adalah frekuensi resonansi praktis (titik tempat $C(\omega)$ maksimal). Kesimpulan ini, misalnya, mengikuti dari uji turunan pertama karena dalam kasus ini $C'(\omega) > 0$ untuk $\omega$ kecil. Sebaliknya, jika $\omega_0^2 - 2p^2$ tidak positif, $C(\omega)$ mencapai maksimum pada $\omega=0$, dan tidak ada resonansi praktis karena dalam sistem kita mengandaikan $\omega > 0$. Dalam kasus ini, amplitudo membesar ketika frekuensi pemaksa mengecil.

Jika resonansi praktis terjadi, frekuensinya lebih kecil daripada $\omega_0$. Ketika redaman $c$ (dan $p$) mengecil, frekuensi resonansi praktis menuju $\omega_0$. Jadi ketika redaman sangat kecil, $\omega_0$ merupakan hampiran yang baik bagi frekuensi resonansi praktis. Perilaku ini sesuai dengan pengamatan bahwa ketika $c=0$, $\omega_0$ adalah frekuensi resonansi.

Pengamatan menarik lain ialah bahwa ketika $\omega \to \infty$, maka $C \to 0$. Artinya, jika frekuensi pemaksa terlalu tinggi, gaya tidak berhasil menggerakkan massa dalam sistem massa-pegas. Secara intuitif ini masuk akal. Jika kita bergerak maju-mundur dengan sangat cepat saat duduk di ayunan, ayunan sama sekali tidak akan bergerak, sekuat apa pun gerakannya. Getaran cepat saling menghapus sebelum massa sempat merespons dengan bergerak ke salah satu arah.

Perilakunya lebih rumit jika fungsi pemaksa bukan gelombang kosinus persis, melainkan misalnya \myindex{gelombang persegi}. Fungsi periodik umum merupakan jumlah (superposisi) banyak gelombang kosinus dengan frekuensi berbeda. Pembaca dianjurkan kembali ke bagian ini setelah mempelajari deret Fourier.

\subsection{Latihan}

\begin{exercise}
Turunkan rumus $x_{sp}$ jika persamaannya $m x'' + c x' + kx = F_0 \sin (\omega t)$. Andaikan $c > 0$.
\end{exercise}

\begin{exercise}
Turunkan rumus $x_{sp}$ jika persamaannya $m x'' + c x' + kx = F_0 \cos (\omega t) + F_1 \cos (3\omega t)$. Andaikan $c > 0$.
\end{exercise}

\begin{exercise}
Ambil $m x'' + c x' + kx = F_0 \cos (\omega t)$. Tetapkan $m > 0$, $k > 0$, dan $F_0 > 0$. Perhatikan fungsi $C(\omega)$. Untuk nilai $c$ berapa (nyatakan dalam $m$, $k$, dan $F_0$) tidak terjadi resonansi praktis (yakni, untuk nilai $c$ berapa tidak ada maksimum $C(\omega)$ untuk $\omega > 0$)?
\end{exercise}

\begin{exercise}
Ambil $m x'' + c x' + kx = F_0 \cos (\omega t)$. Tetapkan $c > 0$, $k > 0$, dan $F_0 > 0$. Perhatikan fungsi $C(\omega)$. Untuk nilai $m$ berapa (nyatakan dalam $c$, $k$, dan $F_0$) tidak terjadi resonansi praktis (yakni, untuk nilai $m$ berapa tidak ada maksimum $C(\omega)$ untuk $\omega > 0$)?
\end{exercise}

\begin{exercise}
\pagebreak[3]
Menara air saat gempa bumi bertindak sebagai sistem massa-pegas. Andaikan wadah di atas penuh dan air tidak bergerak-gerak. Wadah bertindak sebagai massa dan penyangga sebagai pegas, dengan getaran terinduksi yang horizontal. Wadah berisi air mempunyai massa $m=\unit[10,000]{kg}$. Diperlukan gaya 1000 newton untuk memindahkan wadah 1 meter. Demi kesederhanaan, andaikan tanpa gesekan. Saat gempa melanda, menara air dalam keadaan diam (tidak bergerak).
%
Gempa menginduksi gaya eksternal $F(t) = m A \omega^2 \cos (\omega t)$ newton.
\begin{tasks}
\task Berapa frekuensi alami menara air?
\task Jika $\omega$ bukan frekuensi alami, temukan rumus amplitudo maksimum osilasi wadah air (simpangan maksimum dari posisi diam). Geraknya merupakan gelombang berfrekuensi tinggi yang dimodulasi gelombang berfrekuensi rendah, jadi cukup cari konstanta di depan sinus.
\task Andaikan $A = 1$ dan datang gempa dengan frekuensi 0.5 siklus per detik. Berapa amplitudo osilasinya? Andaikan menara runtuh jika wadah air bergerak lebih dari 1.5 meter dari posisi diam. Apakah menara akan runtuh?
\end{tasks}
\end{exercise}

\setcounter{exercise}{100}

\begin{exercise}
Massa \unit[4]{kg} pada pegas dengan $k=\unitfrac[4]{N}{m}$ dan konstanta redaman $c=\unitfrac[1]{Ns}{m}$. Andaikan $F_0 = \unit[2]{N}$. Dengan fungsi pemaksa $F_0 \cos (\omega t)$, temukan $\omega$ yang menyebabkan resonansi praktis dan amplitudo resonansi praktisnya.
\end{exercise}
\exsol{%
$\omega = \frac{\sqrt{31}}{4\sqrt{2}} \approx \unitfrac[0.984]{rad}{s}$ \quad
$C(\omega) = \frac{16}{3\sqrt{7}} \approx \unit[2.016]{m}$
}

\begin{exercise}
Turunkan rumus $x_{sp}$ untuk $mx''+cx'+kx = F_0 \cos(\omega t) + A$, dengan $A$ suatu konstanta. Andaikan $c > 0$.
\end{exercise}
\exsol{%
$x_{sp} =
\frac{(\omega_0^2-\omega^2) F_0}
{m{(2\omega p)}^2+m{(\omega_0^2-\omega^2)}^2} \cos (\omega t) +
\frac{2 \omega p F_0}
{m{(2\omega p)}^2+m{(\omega_0^2-\omega^2)}^2} \sin (\omega t)
+ \frac{A}{k}$,
di mana
$p = \frac{c}{2m}$ dan $\omega_0 = \sqrt{\frac{k}{m}}$.
}

\begin{exercise}
Andaikan tidak ada redaman dalam sistem massa dan pegas dengan $m = \unit[5]{kg}$, $k= \unitfrac[20]{N}{m}$, dan $F_0 = \unit[5]{N}$. Andaikan $\omega$ dipilih tepat sebagai frekuensi resonansi.
\begin{tasks}
\task Temukan $\omega$.
\task Temukan amplitudo osilasi pada waktu $t=100$ detik, jika sistem diam pada $t=0$.
\end{tasks}
\end{exercise}
\exsol{%
a) $\omega = \unitfrac[2]{rad}{s}$ \quad
b) $\unit[25]{m}$
}

